\documentclass[12pt,a4paper]{article}

\usepackage{amsmath,amssymb,amsthm}
\usepackage{mathtools}
\usepackage{booktabs}
\usepackage{array}
\usepackage{hyperref}
\usepackage[margin=1in]{geometry}
\usepackage{enumitem}
\usepackage{graphicx}
\usepackage{xcolor}

\newtheorem{theorem}{Theorem}
\newtheorem{proposition}[theorem]{Proposition}
\newtheorem{lemma}[theorem]{Lemma}
\newtheorem{corollary}[theorem]{Corollary}
\newtheorem{conjecture}[theorem]{Conjecture}
\newtheorem{definition}[theorem]{Definition}
\newtheorem{observation}[theorem]{Observation}
\newtheorem{remark}[theorem]{Remark}

\newcommand{\R}{\mathbb{R}}
\newcommand{\C}{\mathbb{C}}
\newcommand{\Z}{\mathbb{Z}}
\newcommand{\Q}{\mathbb{Q}}
\newcommand{\N}{\mathbb{N}}
\newcommand{\ket}[1]{|#1\rangle}
\newcommand{\bra}[1]{\langle#1|}
\newcommand{\braket}[2]{\langle#1|#2\rangle}
\newcommand{\KS}{\mathrm{KS}}
\newcommand{\CK}{\mathrm{CK}\text{-}31}

\title{The Algebraic Landscape of Kochen--Specker Sets\\in Dimension Three}

\author{Michael Kernaghan}

\date{\today}

\begin{document}

\maketitle

\begin{abstract}
We present a computational survey of Kochen--Specker (KS) uncolorability in three-dimensional Hilbert space across two-symbol coordinate alphabets $\mathcal{A} = \{0, \pm 1, \pm x\}$ drawn from quadratic, cyclotomic, and golden-ratio number fields. In every tested alphabet, KS sets arise only when $x$ supports one of two cancellation mechanisms: \emph{modulus-2 cancellation} (the generator satisfies $|x|^2 = 2$, as in $|\sqrt{2}|^2=2$, $|\sqrt{-2}|^2=2$, or $|\alpha|^2=2$; the integer case $1+1=2$ is the degenerate additive instance) or \emph{phase cancellation} (a vanishing sum of unit-modulus terms, as in $1+\omega+\omega^2=0$). Alphabets whose generators have $|x|^2 \geq 3$ and are not roots of unity produce orthogonal triples but not KS-uncolorability in our survey. This empirical pattern explains why constructions cluster into at least six discrete algebraic islands among the tested fields (with a seventh, cubic island confirmed at higher cost). Two yield potentially new KS graph types: the Heegner-7 ring $\Z[(1+\sqrt{-7})/2]$ (43~vectors) and the golden ratio field $\Q(\varphi)$ (52~vectors, revealed only by cross-product completion); $\Z[\sqrt{-2}]$ provides a new algebraic realization of a known Peres-type graph. Using SAT-based bipartite KS-uncolorability, we verify and extend the input counts of Trandafir and Cabello for bipartite perfect quantum strategies across all six islands. Whether the two-mechanism pattern extends to all number fields remains an open question.
\end{abstract}

\section{Introduction}

The Kochen--Specker (KS) theorem~\cite{KochenSpecker1967} is one of the foundational no-go results of quantum mechanics, demonstrating that quantum observables cannot be assigned pre-existing definite values independently of the measurement context (for a comprehensive review, see Ref.~\cite{Budroni2022}). The theorem is typically proved by exhibiting a finite set of rays (one-dimensional subspaces) in a Hilbert space~$\mathcal{H}$ such that no consistent $\{0,1\}$-valued assignment exists satisfying the constraints imposed by~orthogonality.  While KS proofs exist in any dimension $d \geq 3$---including constructions with 20~vectors in $d = 4$~\cite{Kernaghan1994} and 36~vectors in $d = 8$~\cite{KernaghanPeres1995}---dimension~3 remains the most intensively studied case because it admits the smallest known sets and the tightest algebraic constraints.

In dimension $d = 3$, the smallest known KS set is the Conway--Kochen 31-vector set ($\CK$), reported by Peres~\cite{Peres1993} with coordinates drawn from the integer alphabet $\{0, \pm 1, \pm 2\}$. Despite decades of effort, no smaller construction has been found. The best known lower bound is 24 vectors, established independently by Li, Bright, and Ganesh~\cite{LiBrightGanesh2024} and by Kirchweger, Peitl, and Szeider~\cite{KirchwegarPeitlSzeider2023} using SAT-based methods with algebraic realizability checking (improving the earlier bound of 22 by Uijlen and Westerbaan~\cite{UijlenWesterbaan2016}), leaving a gap of 7 vectors that remains one of the central open problems in KS theory. Trandafir and Cabello~\cite{TrandafirCabello2025} recently proved that $\CK$ is rigid (unique up to unitary transformation in $\C^3$) and conjectured that 31 is optimal.

Previous computational approaches have primarily searched within fixed coordinate alphabets~\cite{Peres1991, Pavicic2004} or generated random abstract hypergraphs and tested realizability~\cite{LiBrightGanesh2024}. The algebraic structure of the coordinate systems underlying KS sets---specifically, which number fields can support them---has received less systematic attention, though Cortez, Morales, and Reyes~\cite{CortezMoralesReyes2022} showed (in a preprint whose arguments we have verified) that $M_3(\Z[1/6])_{\mathrm{sym}}$---the partial ring of $3 \times 3$ symmetric matrices over $\Z[1/6]$---is the minimal ring extension of $\Z$ admitting no algebraic hidden states, a ring-theoretic formulation of KS contextuality in dimension~3. (Indeed, the original Kochen--Specker construction~\cite{KochenSpecker1967} uses coordinates from the real cyclotomic field $\mathbb{Q}(\cos\pi/10)$, with the 18-degree rotations that chain its 15 copies of the fundamental 10-ray gadget introducing algebraic irrationals related to the golden ratio---an early, if implicit, instance of the alphabet-dependence that the present paper makes explicit.)

In this paper, we ask a structural question: \emph{what algebraic property of a coordinate alphabet controls whether it can support KS sets in dimension 3?} Through a systematic computational survey of quadratic fields $\Q(\sqrt{d})$ for $d = 2, \ldots, 30$, all nine imaginary quadratic fields with class number~1, cyclotomic fields, and golden ratio extensions, we find a consistent empirical pattern: in every tested case, KS-uncolorability requires the alphabet to support a \emph{cancellation identity}---an algebraic relation among its elements that makes dot products vanish exactly---of one of two types: \emph{modulus-2 cancellation} (the generator $x$ satisfies $|x|^2 = 2$) or \emph{phase cancellation} (a vanishing sum of unit-modulus terms). Among all tested alphabets, these are the only two mechanisms observed to produce the dense triad networks required for KS-uncolorability, explaining why constructions cluster into at least six discrete ``algebraic islands'' among the tested fields (a seventh, cubic island is confirmed in Section~\ref{sec:limitations}). Whether this pattern extends to all number fields remains an open question.

The two-mechanism pattern unifies all known constructions: the integer identity $1 + 1 = 2$, the Peres identity $|\sqrt{2}|^2 = 2$, the Eisenstein identity $1 + \omega + \omega^2 = 0$, and the Heegner-7 identity $|\alpha|^2 = 2$ are all instances of modulus-2 or phase cancellation.  Alphabets whose generators have $|x|^2 \geq 3$ and are not roots of unity produce orthogonal pairs and triples but never enough interlocking triads for KS-uncolorability in our survey.  Our main contributions are:

\begin{enumerate}[label=(\roman*)]
    \item An empirical classification of the cancellation identities that support KS-un\-color\-ability in two-symbol alphabets $\{0, \pm 1, \pm x\}$: across all tested quadratic, cyclotomic, and golden-ratio fields, KS sets arise only from modulus-2 or phase cancellation, with no intermediate mechanisms observed (Sections~\ref{sec:quadratic}--\ref{sec:heegner}).
    \item Two KS configurations that, to our knowledge, have not been previously reported: a 43-vector set from the Heegner-7 ring $\Z[\alpha]$ with $\alpha = (1{+}\sqrt{-7})/2$ (whose orthogonality graph does not match any known KS construction in published catalogues we have checked), and a 52-vector set from the golden ratio field $\Q(\varphi)$, revealed only by cross-product completion (Sections~\ref{sec:heegner},~\ref{sec:golden}; Appendix~\ref{app:rays}).
    \item Evidence that the algebraic substrate has operational consequences: different islands yield different bipartite perfect quantum strategies (Section~\ref{sec:cabello}) and different Cabello--Severini--Winter contextual advantages (Section~\ref{sec:csw}), with no single island dominating on all measures.
    \item Supporting results: an OCUS-certified proof that no $\leq 30$-ray KS subset exists within the full 49-ray integer pool (Section~\ref{sec:exhaustive}), universal merge saturation across all six islands (Observation~\ref{obs:merge}), a Jacobian-based rigidity analysis confirming that four of six minimal KS sets are rigid in~$\C^3$ (Observation~\ref{obs:rigidity}), as well as identifications of three distinct 33-vector KS sets with pairwise non-isomorphic graphs (Observation~\ref{obs:three-33}), a MUS landscape analysis with a norm-stratified invariant core (Observation~\ref{obs:mus-landscape}), the $6 \mid n$ cyclotomic theorem (Theorem~\ref{thm:6n}), and graph-isomorphism evidence for the modulus-2 family (Proposition~\ref{prop:isomorphism}).
\end{enumerate}

\section{Preliminaries}\label{sec:prelim}

\subsection{Kochen--Specker sets and colorings}

\begin{definition}
A \emph{ray} in $\mathcal{H} = \C^d$ (or $\R^d$) is a one-dimensional subspace, represented by a nonzero vector up to scalar multiples. Two rays $v, w$ are \emph{orthogonal} if $\braket{v}{w} = 0$. An \emph{orthogonal triad} in $d = 3$ is a set of three mutually orthogonal rays forming a complete basis.
\end{definition}

\begin{definition}
A \emph{KS coloring} of a set of rays $S$ is a function $f: S \to \{0, 1\}$ (interpreted as ``red'' and ``green'') satisfying:
\begin{enumerate}[label=(\alph*)]
    \item \emph{Orthogonal exclusion}: if $v \perp w$ and $f(v) = 1$, then $f(w) = 0$;
    \item \emph{Triad completeness}: for every orthogonal triad $\{v_1, v_2, v_3\} \subset S$, exactly one $v_i$ has $f(v_i) = 1$.
\end{enumerate}
A set $S$ is a \emph{Kochen--Specker set} (or \emph{KS-uncolorable}) if no KS coloring exists.
\end{definition}

\begin{definition}
The \emph{orthogonality graph} $G(S)$ of a ray set $S$ has vertices corresponding to rays and edges connecting orthogonal pairs. A \emph{triad} corresponds to a triangle (3-clique) in $G(S)$. We use ``orthogonality graph'' when referring to the pair structure and ``orthogonality hypergraph'' when emphasizing the triad (3-uniform hyperedge) structure; both encode the same underlying ray configuration.
\end{definition}

\subsection{Coordinate alphabets and number fields}

A \emph{coordinate alphabet} $\mathcal{A} \subset \R$ (or $\C$) is a finite set of values from which ray coordinates are drawn.  (Pavicic et al.~\cite{Pavicic2019} call these ``vector components''; we adopt ``alphabet'' to emphasize the combinatorial generating role---one selects symbols from $\mathcal{A}$ and assembles them into three-coordinate ``words'' that represent rays.)  Given an alphabet $\mathcal{A}$, the induced ray set is
\[
S(\mathcal{A}) = \{ [v] : v \in \mathcal{A}^3 \setminus \{0\} \}
\]
where $[v]$ denotes the equivalence class of $v$ under scalar multiplication (the ray).  The alphabet of a known KS set can often be read directly from its geometric representation: in the cube diagrams used to display the Peres~33, Conway--Kochen~31, and related constructions~\cite{Peres1991,Peres1993}, each vertex carries explicit integer coordinates, and the set of distinct coordinate values appearing across all vertices is precisely the alphabet $\mathcal{A}$.

Two rays $v, w \in \mathcal{A}^3$ are orthogonal when their dot product vanishes: $\sum_k \bar v_k w_k = 0$.  For coordinates drawn from a finite alphabet, this requires an algebraic relation among elements of $\mathcal{A}$ that makes the sum cancel exactly---a \emph{cancellation identity}.  The cancellation identity is the mechanism by which an alphabet produces orthogonal rays, and hence the mechanism by which it can potentially generate a Kochen--Specker set.

For the integer alphabet $\{0, \pm 1, \pm 2\}$, the key identity is $1 + 1 = 2$: two coordinates of magnitude~1 can sum to cancel a single coordinate of magnitude~2, as in $1 \cdot 1 + 1 \cdot 1 + (-1) \cdot 2 = 0$.  For the Peres alphabet $\{0, \pm 1, \pm\sqrt{2}\}$, it is $(\sqrt{2})^2 = 2 = 1 + 1$.  For the Eisenstein alphabet $\{0, \pm 1, \pm\omega, \pm\bar\omega\}$ (where $\omega = e^{2\pi i/3}$), it is the three-term phase cancellation $1 + \omega + \omega^2 = 0$.  In every case, the identity has the same effect: it creates exact orthogonalities among triples of rays (triads), and when enough triads interlock, no consistent $\{0,1\}$-coloring exists---a KS set.

One might ask whether the trivial identity $1 + (-1) = 0$ already suffices---after all, it makes dot products vanish (e.g., $1 \cdot 1 + 1 \cdot (-1) + 0 \cdot 0 = 0$).  The answer is that the minimal alphabet $\{0, \pm 1\}$ generates only 13 projective rays in $\R^3$ (the 3 axis directions, 6 face diagonals, and 4 body diagonals of the cube), far too few for a KS set (the minimum known is 31).  The identity $1 + 1 = 2$ is what unlocks the $\pm 2$ coordinate, expanding the pool from 13 to 49 rays---nearly four times as many---and creating the dense triad network that KS-uncolorability requires.  In this sense, $1 + 1 = 2$ is the simplest \emph{useful} cancellation identity: the first one that generates enough orthogonal structure for the coloring constraints to become unsatisfiable.

We classify the algebraic mechanisms that produce orthogonalities.

\begin{definition}\label{def:cancellation}
Let $\mathcal{A} = \{0, \pm 1, \pm x\}$ be a two-symbol coordinate alphabet.  The \emph{product set} $\bar{\mathcal{A}} \cdot \mathcal{A} = \{\bar{a}b : a, b \in \mathcal{A}\}$ contains all values that can appear as individual terms in a dot product $\langle v | w \rangle = \sum_k \bar{v}_k w_k$.  A \emph{cancellation identity} is a vanishing sum $\sum_{k} t_k = 0$ with $t_k \in \bar{\mathcal{A}} \cdot \mathcal{A}$ (at most three terms for $d = 3$).  We distinguish two mechanisms:
\begin{itemize}
\item \emph{Modulus-2 cancellation}: the generator $x$ satisfies $|x|^2 = 2$ (Hermitian squared modulus), as in $|\sqrt{2}|^2 = 2$, $|\sqrt{-2}|^2 = 2$, $|1+i|^2 = 2$, or $|\alpha|^2 = 2$ for $\alpha = (1+\sqrt{-7})/2$.  The integer case $1+1=2$ is the degenerate instance where $x = 2$ itself and the identity arises additively.
\item \emph{Phase cancellation}: $x$ is a root of unity and a vanishing sum of unit-modulus terms exists, as in $1 + \omega + \omega^2 = 0$ for $\omega = e^{2\pi i/3}$.
\end{itemize}
We say the alphabet has \emph{low-complexity cancellation} if it supports one of these mechanisms.  Note that a cancellation identity is a \emph{local} condition enabling individual orthogonalities; KS-uncolorability additionally requires enough interlocking triads for the coloring constraints to be globally unsatisfiable.  In the cyclotomic case, phase cancellation ($3 \mid n$) produces triads but not KS-uncolorability; the additional condition $2 \mid n$ is needed for sufficient interlocking (Theorem~\ref{thm:6n}).
\end{definition}

\begin{remark}[Terminology]\label{rem:norm-clarification}
We use ``modulus-2'' to refer to the Hermitian squared modulus $|x|^2$ of the generator, \emph{not} to the algebraic field norm $N_{K/\Q}(x)$ (which can differ: $N_{\Q(\sqrt{5})/\Q}(\varphi) = -1$, not~$\varphi^2$) nor to the maximum squared modulus of individual dot-product terms (which may exceed~2 even in the integer case: the cancellation $1 + 1 - 2 = 0$ has $\max|t_k|^2 = 4$).  The salient property is that the generator produces the value~2 through its squared modulus, enabling three-term cancellations with bounded coordinate size.
\end{remark}

For the six known KS-supporting alphabets, the generating relations are: integer ($1+1=2$, modulus-2), Peres ($|\sqrt{2}|^2=2$, modulus-2), Eisenstein ($1+\omega+\omega^2=0$, phase), $\Z[\sqrt{-2}]$ ($|\sqrt{-2}|^2=2$, modulus-2), Heegner-7 ($|\alpha|^2=2$, modulus-2), and golden ($\varphi^2=\varphi+1$---the raw alphabet lacks low-complexity cancellation, but cross-product completion introduces $1/\varphi$ with $|1/\varphi|^2 < 1$, creating effective modulus-2 cancellations in the completed pool; see Section~\ref{sec:golden}).  No alphabet whose generator has $|x|^2 \geq 3$ and is not a root of unity produces a KS set in our survey.

The methodological approach of this paper is to \emph{reverse} this chain: we systematically construct coordinate alphabets from algebraic number fields, identify their cancellation identities, generate the induced ray sets, and test for KS-uncolorability.  This turns the search for new KS sets into a question of algebra: which rings admit cancellation identities of low enough norm to produce the dense orthogonality structures that KS-uncolorability requires?

\subsection{Cross-product completion}

\begin{definition}\label{def:completion}
The \emph{cross-product completion} of a ray set $S \subset \R^3$ is the closure $\bar{S}$ obtained by iterating: for every orthogonal pair $v, w \in S$, add the ray $[v \times w]$ to $S$, until no new rays are generated. This is well-defined since $v \perp w$ implies $v \times w \neq 0$ and $v \times w$ is orthogonal to both $v$ and $w$. Termination is not guaranteed in general but is observed empirically for all alphabets in our survey, typically stabilizing within 3--5 iterations (see Section~\ref{sec:limitations}).
\end{definition}

Cross-product completion is a natural choice of closure operation: it ensures that every orthogonal pair is extended to a complete basis, which is the minimal requirement for a ray set to carry KS-type coloring constraints.  Other closure operations are possible (e.g., closure under orthogonal complementation of arbitrary subsets, or Gram--Schmidt processes with field constraints), and the island classification may in principle depend on which closure is used.  We adopt cross-product completion throughout and note that it terminates empirically for all tested alphabets; a proof of termination for specific families of coordinate rings would strengthen the classification.  Since cross products are polynomial in the coordinates, completion trivially preserves any number field $K$. The nontrivial property is \emph{finiteness}: for a given alphabet $\mathcal{A}$, the completed ray pool $\bar{S}(\mathcal{A})$ stabilizes after finitely many iterations, and its size and triad structure depend sensitively on the algebraic identities available in $\mathcal{A}$.

\emph{Important caveat}: completion can introduce coordinate values absent from the original alphabet, thereby creating cancellation identities that $\mathcal{A}$ alone does not support.  The golden ratio island (Section~\ref{sec:golden}) is a concrete example: the raw alphabet $\{0, \pm 1, \pm\varphi\}$ produces a colorable set, but completion generates the reciprocal $1/\varphi$, introducing effective modulus-2 cancellations that make the completed pool KS-uncolorable.  Consequently, the island classification in this paper is properly a classification of \emph{completed coordinate algebras}---the closure of the alphabet under cross products---rather than of the initial alphabets alone.  The raw alphabet serves as a convenient generating set, but the algebraic structure relevant to KS-uncolorability is that of the completed pool.

\section{Computational Methods}\label{sec:methods}

\subsection{Ray canonicalization}\label{sec:canon}

A ray is an equivalence class of nonzero vectors under scalar multiplication. For \emph{real} alphabets ($\mathcal{A} \subset \R$), two vectors represent the same ray iff they are nonzero real multiples of each other; we choose the primitive representative with first nonzero entry positive. For \emph{complex} alphabets ($\mathcal{A} \subset \C$), two vectors represent the same ray iff they are nonzero complex multiples of each other (i.e., $v \sim \lambda v$ for any $\lambda \in \C^*$); we choose a canonical representative by dividing by the first nonzero coordinate (making it equal to 1), then applying a lexicographic tie-breaker on the remaining coordinates.

In both cases, representatives are \emph{unnormalized}: we do not divide by the vector's norm. This avoids introducing spurious irrationals through normalization and ensures that ray counts (e.g., 49, 109, 205) are well-defined projective counts over the coordinate ring. The canonicalization is deterministic and injective on rays: each projective equivalence class maps to exactly one representative. For real alphabets, the sign convention (first nonzero entry positive) resolves the $\pm v$ ambiguity; for complex alphabets, dividing by the first nonzero coordinate resolves the full $\C^*$ scaling ambiguity. Orthogonality is tested on the unnormalized representatives: $v \perp w$ iff $\sum_k \bar{v}_k w_k = 0$ (Hermitian, reducing to $\sum_k v_k w_k = 0$ in the real case). Since orthogonality is invariant under scaling, this test is independent of the choice of representative.

\smallskip\noindent\textbf{Exact arithmetic.}  All coordinates are represented as exact algebraic numbers: Python integers for the integer and Peres alphabets, and tuples $(a, b) \mapsto a + b\sqrt{d}$ (or $(a,b) \mapsto a + b\omega$) for quadratic and Eisenstein fields, with conjugation and multiplication implemented symbolically.  For higher cyclotomic fields ($n > 6$), coordinates are roots of unity $\zeta_n^k$ with exact multiplication via exponent arithmetic modulo~$n$ and exact conjugation $\overline{\zeta_n^k} = \zeta_n^{n-k}$; the $6 \mid n$ theorem is proved algebraically rather than by exhaustive coordinate computation.  For real irrational alphabets (golden ratio, cubic), coordinates are represented as tuples over the minimal polynomial basis (e.g., $(a,b) \mapsto a + b\varphi$ for $\Q(\varphi)$) with exact multiplication.  Orthogonality is certified by verifying that the Hermitian dot product vanishes \emph{exactly} (integer or algebraic zero), with no floating-point approximation.  This guarantees that all reported orthogonality relations, triad counts, and SAT encodings are exact.

\subsection{KS coloring via SAT}

We encode the KS coloring problem as a Boolean satisfiability (SAT) instance. For a ray set with $n$ rays, orthogonal pairs $E$, and triads $T$:

\begin{itemize}
    \item Boolean variable $x_i$ indicates ray $i$ is green ($f(i) = 1$).
    \item For each triad $\{i, j, k\} \in T$: exactly-one constraint (4 clauses).
    \item For each pair $(i, j) \in E$: at-most-one constraint (1 clause: $\neg x_i \lor \neg x_j$).
\end{itemize}

We use the Glucose4 solver~\cite{AudemardSimon2018} via the PySAT library. UNSAT = KS-uncolorable. For a typical 33-ray, 16-triad configuration, the SAT instance has 33 variables and $4 \times 16 + |E| \approx 120$--$185$ clauses (depending on the number of orthogonal pairs $|E|$). No symmetry breaking is applied; the instances are small enough for direct solving (milliseconds). All SAT instances, ray coordinates, and basis lists are available in the accompanying code repository.

\subsection{Randomized greedy minimization}

Given a KS-uncolorable set $S$, we find minimal KS subsets by randomized greedy removal: shuffle the rays, attempt to remove each in random order, keeping the removal if the set remains KS-uncolorable. We run 200--1000 independent trials per configuration.  The \emph{pool minimum} is the smallest KS subset size within a given finite ray pool; for each island, we certify this minimum using OCUS (Section~\ref{sec:exhaustive}).

\subsection{Realizability via numerical optimization}

Given an abstract orthogonality graph $G$ on $n$ vertices with edge set $E$, we test realizability in $\R^3$ by parametrizing each ray by spherical coordinates $(\theta_i, \phi_i)$ and minimizing
\[
\mathcal{L}(\theta, \phi) = \sum_{(i,j) \in E} (v_i \cdot v_j)^2
\]
using L-BFGS-B with 50 random restarts. We declare realizability if $\mathcal{L} < 10^{-8}$.

\section{Real Coordinate Alphabets}\label{sec:real}

\subsection{Systematic alphabet survey}

Table~\ref{tab:alphabets} summarizes the results of exhaustive ray generation and KS testing for coordinate alphabets based on various algebraic extensions of $\Q$.

\begin{table}[ht]
\centering
\caption{KS sets from real coordinate alphabets $\{0, \pm 1, \pm x\}$ in $\R^3$.}
\label{tab:alphabets}
\begin{tabular}{lcccccc}
\toprule
Alphabet & Rays & Pairs & Triads & KS? & Min \\
\midrule
$\{0, \pm 1\}$ & 13 & 24 & 4 & No & --- \\
$\{0, \pm 1, \pm\sqrt{2}\}$ & 49 & 120 & 16 & \textbf{Yes} & \textbf{33} \\
$\{0, \pm 1, \pm 2\}$ & 49 & 138 & 26 & \textbf{Yes} & \textbf{31} \\
$\{0, \pm 1, \pm\sqrt{3}\}$ & 49 & 114 & 10 & No & --- \\
$\{0, \pm 1, \pm\sqrt{5}\}$ & 49 & 114 & 10 & No & --- \\
$\{0, \pm 1, \pm\varphi\}$ & 49 & 138 & 10 & No & --- \\
$\{0, \pm 1, \pm\frac{1}{2}\}$ & 49 & 138 & 26 & \textbf{Yes} & \textbf{31} \\
\bottomrule
\end{tabular}
\end{table}

The Conway--Kochen 31-vector set is rediscovered by the integer alphabet in 19\% of randomized minimization trials (192/1000). No trial produced a set smaller than 31.

\subsection{The 2:1 ratio as the critical feature}\label{sec:ratio}

\begin{observation}
Every real alphabet that achieves a 31-vector KS set contains the ratio $2{:}1$ among its nonzero elements---either as $\pm 2$ alongside $\pm 1$, or equivalently as $\pm \frac{1}{2}$ alongside $\pm 1$. Adding other algebraic values (e.g., $\sqrt{2}$, $\varphi$, $\sqrt{3}$) to such an alphabet does not reduce the minimum below 31.
\end{observation}

Table~\ref{tab:extended} shows extended alphabets. In every case, the presence of the 2:1 ratio yields 31; its absence yields $\geq 33$ or colorable.

\begin{table}[ht]
\centering
\caption{Extended alphabets: across all tested configurations, the 2:1 ratio correlates perfectly with reaching 31.}
\label{tab:extended}
\begin{tabular}{lcccc}
\toprule
Alphabet & Rays & Triads & Min KS & Has 2:1? \\
\midrule
$\{0,\pm 1,\pm\sqrt{2},\pm 2\}$ & 109 & 38 & 31 & Yes \\
$\{0,\pm 1,\pm\varphi,\pm 2\}$ & 145 & 50 & 31 & Yes \\
$\{0,\pm 1,\pm 2,\pm 3\}$ & 145 & 50 & 31 & Yes \\
$\{0,\pm 1,\pm\sqrt{2},\pm\varphi\}$ & 145 & 28 & 33 & No \\
$\{0,\pm 1,\pm\varphi,\pm\varphi^2\}$ & 109 & 24 & colorable & No \\
\bottomrule
\end{tabular}
\end{table}

Among the two-element alphabets tested, $1 + 1 = 2$ produces the highest triad density: 26 triads from 49 rays. This density appears to be the mechanism by which the integer alphabet achieves the smallest minimum of 31. Heuristically, the identity $a + a = b$ with $a, b$ both in the alphabet generates a large number of three-term dot-product cancellations and their coordinate permutations. Identities involving irrationals (e.g., $(\sqrt{2})^2 = 2$) produce fewer cancellations because they require squaring to reach an integer, reducing the number of distinct triad-generating permutations. We note that ``unique'' here is relative to the alphabets in our survey and the equivalence class under scaling; a formal proof of optimality among all possible two-element cancellation identities remains open.

\subsection{Computational evidence that 31 is minimum for the integer island}\label{sec:exhaustive}

We provide strong computational evidence that no smaller KS set exists within the integer island's ray pool, though a complete proof over the full pool remains computationally infeasible.

Starting from the 49-ray pool of $\{0, \pm 1, \pm 2\}$, we ran 1{,}000 independent randomized greedy minimization trials, each producing a minimal KS subset. Six distinct 31-ray sets were found, collectively using exactly 37 of the 49 available rays. We then performed an exhaustive check: every $\binom{37}{30} = 10{,}295{,}472$ subset of size 30 drawn from these 37 rays was tested for KS-uncolorability via SAT. All were colorable.

\begin{proposition}\label{prop:31optimal}
No 30-ray subset of the 37-ray union of all discovered minimal 31-sets is KS-uncolorable. Furthermore, no single-ray removal, two-ray removal, or ray-swap operation applied to any of the six 31-sets produces a KS-uncolorable 30-ray configuration.
\end{proposition}

\begin{remark}[Optimality within the full pool]\label{rem:ocus}
The exhaustive check above covers the 37-ray union of discovered 31-sets, not the entire 49-ray pool (49 rays, 138 orthogonal pairs, 26 triads).  We close this gap using MUS (Minimal Unsatisfiable Subset) extraction.

\smallskip\noindent\textbf{MUS-based certification.}  The KS coloring formula for the full 49-ray pool---encoding constraints (I) and (II) over all triads and orthogonal pairs---is unsatisfiable (the pool is KS-uncolorable).  We extract ray-level MUSes using assumption-based solving: each ray $r$ has a selector literal $a_r$, and all constraints involving $r$ are conditioned on $a_r$. The solver returns an unsatisfiable core in terms of these selectors, which we then minimize by iterative ray deletion. The result is a minimal subset of rays whose induced constraints are unsatisfiable.  Any MUS therefore identifies a \emph{minimal KS-uncolorable subset} of the pool.  We performed 1{,}000 independent MUS extractions (using randomized deletion with PySAT/Glucose4); all returned subsets of exactly 31~rays, yielding 162 distinct minimal KS sets (6 up to the pool's automorphism group).

\smallskip\noindent\textbf{Exhaustive lower bound.}  An OCUS (Optimal Constrained Unsatisfiable Subset) procedure subsequently proved that no KS-uncolorable subset of $\leq 30$ rays exists within the full 49-ray pool, exhausting all candidates in 272 iterations ($< 0.2$\,s).  This certifies that 31 is the exact minimum within the pool.  As independent confirmation, we tested every $\binom{37}{30}$ subset of the 37-ray union of all discovered minimal sets; all were colorable.
\end{remark}

The exhaustive check in Proposition~\ref{prop:31optimal} required approximately 10.3 million SAT calls. As an independent verification, we also performed neighborhood searches around each of the six 31-sets: all 31 single-ray removals (per set), all 8{,}370 ray-swap operations, and all 465 two-ray removals per set were tested. Every resulting 30-ray configuration was colorable.

The OCUS certification proves that 31 is the exact minimum within the 49-ray integer pool.  Combined with the rigidity theorem of Trandafir and Cabello~\cite{TrandafirCabello2025} (CK-31 unique up to unitary equivalence) and the absence of any sub-31 KS set across all six algebraic islands, this is consistent with---but does not prove---the conjecture that 31 is optimal for all KS sets in $\R^3$.

\subsection{Quadratic number fields}\label{sec:quadratic}

We systematically tested the alphabet $\{0, \pm 1, \pm\sqrt{d}\}$ for every non-square integer $d = 2, \ldots, 30$.

\begin{proposition}\label{prop:quadratic}
Among the quadratic fields $\Q(\sqrt{d})$ for $d \in \{2, 3, 5, 6, 7, \ldots, 30\}$ (excluding perfect squares), only $d = 2$ produces a KS-uncolorable set from the standard two-element alphabet $\{0, \pm 1, \pm\sqrt{d}\}$. All others have $\leq 10$ triads and are colorable.
\end{proposition}

A heuristic explanation: every quadratic alphabet has a cancellation identity---$(\sqrt{d})^2 = d$---but the identity produces dense triad networks only when $d$ equals a sum of alphabet elements, i.e., when the cancellation has \emph{low norm}.  For $d = 2$, we have $(\sqrt{2})^2 = 1 + 1$, giving the Peres cancellation with norm~2.  For $d \geq 3$, no two-term sum from $\{1, \sqrt{d}\}$ equals $d$, so the cancellation has norm~$\geq 3$ and the orthogonality structure is too sparse for KS-uncolorability.  The explanatory variable is therefore the \emph{norm of the cancellation}, not the mere existence of one.  (A proof that no other cancellation pattern can arise for these alphabets would require a more detailed case analysis, which we leave to future work.)

\section{Complex Coordinate Alphabets}\label{sec:complex}

\subsection{Roots of unity and the $6 \mid n$ conjecture}\label{sec:roots}

We generated ray sets from the alphabet $\{0\} \cup \{\zeta^k : k = 0, \ldots, n-1\}$ where $\zeta = e^{2\pi i/n}$, using the Hermitian inner product $\braket{v}{w} = \sum_k \bar{v}_k w_k$ for orthogonality. Table~\ref{tab:roots} shows results for $n = 2, \ldots, 30$.

\begin{table}[ht]
\centering
\caption{KS sets from $n$th roots of unity. Only multiples of 6 (bold) are uncolorable. The ``Best found'' column reports the smallest KS subset found by heuristic minimization; these are upper bounds, not proven minima. For $6 \mid n$, subset containment guarantees $\leq 33$ (see text).}
\label{tab:roots}
\begin{tabular}{ccccccc}
\toprule
$n$ & $6 \mid n$? & Rays & Triads & KS? & Best found & $\leq 33$? \\
\midrule
2--5 & No & 13--43 & 1--7 & No & --- & --- \\
\textbf{6} & \textbf{Yes} & \textbf{57} & \textbf{22} & \textbf{Yes} & \textbf{33} & \textbf{Yes} \\
7--11 & No & 73--157 & 1--28 & No & --- & --- \\
\textbf{12} & \textbf{Yes} & \textbf{183} & \textbf{67} & \textbf{Yes} & \textbf{33} & \textbf{Yes} \\
13--17 & No & 211--343 & 1--76 & No & --- & --- \\
\textbf{18} & \textbf{Yes} & \textbf{381} & \textbf{136} & \textbf{Yes} & \textbf{76} & \textbf{Yes} \\
19--23 & No & 421--601 & 1--148 & No & --- & --- \\
\textbf{24} & \textbf{Yes} & \textbf{651} & \textbf{229} & \textbf{Yes} & --- & \textbf{Yes} \\
25--29 & No & 703--931 & 1--244 & No & --- & --- \\
\textbf{30} & \textbf{Yes} & \textbf{993} & \textbf{346} & \textbf{Yes} & --- & \textbf{Yes} \\
\bottomrule
\end{tabular}
\end{table}

\begin{theorem}\label{thm:6n}
The ray set generated by $n$th roots of unity in $\C^3$ is KS-uncolorable if and only if $6 \mid n$.
\end{theorem}

\begin{proof}
\emph{Sufficiency.} If $6 \mid n$, then $\langle\zeta^{n/6}\rangle$ generates the 6th roots of unity as a subgroup; the Eisenstein 33-vector KS set~\cite{Cabello2025simplest} lives in the $n=6$ pool, which embeds into the $n$-pool.  (This is a proof of existence by embedding a known construction; an intrinsic characterization of \emph{why} $6 \mid n$ suffices---rather than inheriting from the Eisenstein set---remains open.)

\emph{Necessity.} We show that if $6 \nmid n$, then the pool is KS-colorable.

\textbf{Vanishing-sum lemma.}  Since inner products in $\C^3$ involve at most three nonzero terms, only 1-, 2-, and 3-term vanishing sums of $n$th roots arise as orthogonality conditions (longer vanishing sums exist for general $n$; see Lam and Leung~\cite{LamLeung2000} for the complete classification, which is not needed here).  A three-term sum $\zeta^a + \zeta^b + \zeta^c = 0$ with $\zeta = e^{2\pi i/n}$ exists if and only if $3 \mid n$.  \emph{Proof}: dividing by $\zeta^a$ gives $1 + \zeta^p + \zeta^q = 0$; setting $\alpha = \zeta^p$ requires $|-1-\alpha| = 1$, which forces $\cos(2\pi p/n) = -1/2$, hence $p/n \in \{1/3, 2/3\}$, so $3 \mid n$.  Similarly, $\zeta^a + \zeta^b = 0$ requires $\zeta^{b-a} = -1$, hence $2 \mid n$.

\textbf{Case 1: $3 \nmid n$, $2 \nmid n$.}  No 2-term or 3-term vanishing sums of $n$th roots exist, so no dot product of the form $\sum_{k=1}^{3} \bar{v}_k w_k$ (with $v_k, w_k \in \{0\} \cup \langle\zeta\rangle$) can vanish unless both rays share a zero coordinate pattern that eliminates all terms.  The only orthogonalities are axis-complementary, the sole triad is $\{[1{:}0{:}0],\, [0{:}1{:}0],\, [0{:}0{:}1]\}$, and the pool is trivially colorable.

\textbf{Case 2: $3 \nmid n$, $2 \mid n$.}  Two-term cancellations $\zeta^a + \zeta^b = 0$ exist (via $\zeta^{n/2} = -1$), but no three-term vanishing sum exists.  Consequently, no pair of all-nonzero rays is orthogonal; all triads have the form $\{v_1, v_2, e_k\}$ where $v_1, v_2$ are one-zero rays in the same coordinate plane and $e_k$ is the corresponding axis ray.  Within each plane (say $z = 0$), the rays are parameterized by $p = b - a \bmod n$ via $(1, \zeta^p, 0)$, and two rays $p, q$ are orthogonal iff $q - p \equiv n/2 \pmod{n}$---a \emph{perfect matching} (each ray has exactly one partner).  Moreover, one-zero rays from different planes are never orthogonal: $\langle(\zeta^a, \zeta^b, 0) | (\zeta^c, 0, \zeta^d)\rangle = \zeta^{c-a} \neq 0$.  An explicit coloring is: color one axis ray green (forcing its plane's rays all red), color the other two axis rays red (their planes decompose into independent matched pairs, each 2-colorable).

\textbf{Case 3: $3 \mid n$, $2 \nmid n$.}  This is the principal case.  Let $\omega = \zeta^{n/3}$, a primitive cube root of unity.  We prove three lemmas establishing that every triad is isolated.

\emph{Projective collapse lemma.}  For any all-nonzero ray $v = (\zeta^a, \zeta^b, \zeta^c)$, the orthogonality condition $\langle v | w \rangle = 0$ with $w$ all-nonzero requires a three-term vanishing sum, forcing the coordinate differences to be a permutation of $(0, n/3, 2n/3)$.  The six permutations yield six candidate neighbors, but these collapse to exactly \emph{two} projectively distinct rays: the even permutations give $w_+ = [\zeta^a : \omega\zeta^b : \omega^2\zeta^c]$ (since $(1{,}2{,}0) \mapsto \omega \cdot (0{,}1{,}2)$ and $(2{,}0{,}1) \mapsto \omega^2 \cdot (0{,}1{,}2)$), and the odd permutations give $w_- = [\zeta^a : \omega^2\zeta^b : \omega\zeta^c]$.  Moreover, $\langle w_+ | w_- \rangle = 1 + \bar\omega\omega^2 + \bar\omega^2\omega = 1 + \omega + \omega^2 = 0$, so $\{v, w_+, w_-\}$ is a triad---the \emph{unique} triad containing $v$ among all-nonzero rays.

\emph{One-zero isolation lemma.}  For odd $n$, no ray with exactly one zero coordinate participates in any triad.  \emph{Proof}: a triad involving $({\zeta^a}, {\zeta^b}, 0)$ would require a second ray $u$ with $u_3 = 0$ (forced by orthogonality to the complementary axis ray) and $\langle v | u \rangle = \zeta^{e-a} + \zeta^{f-b} = 0$, i.e., $-1 \in \langle\zeta\rangle$, which requires $2 \mid n$.

\emph{Decomposition.}  For odd $n$ with $3 \mid n$, the triads are: one axis triad $\{[1{:}0{:}0],\, [0{:}1{:}0],\, [0{:}0{:}1]\}$ plus $n^2/3$ all-nonzero triads, each isolated (no shared rays).  Total: $n^2/3 + 1$ triads with maximum membership~1 per ray.  The coloring problem decomposes into $n^2/3 + 1$ independent subproblems, each trivially satisfiable.
\end{proof}

\begin{remark}[Subset containment]
For any $n$ divisible by 6, the $n$th roots of unity contain the 6th roots as a subset, so the $n=6$ ray pool embeds into the $n$-ray pool. Consequently, any $6\mid n$ pool contains a 33-vector KS subset (the Eisenstein-based set from $n=6$). The ``Best found'' column in Table~\ref{tab:roots} reports the output of heuristic greedy minimization, which for $n=18$ returned a 76-vector set---a failure of the minimizer, not a true lower bound. The true minimum for every $6\mid n$ pool is at most 33.
\end{remark}

The mechanism requires two ingredients working in concert:
\begin{itemize}
    \item \emph{Three-term phase cancellation} from $3 \mid n$: the identity $1 + \omega + \omega^2 = 0$ (where $\omega = e^{2\pi i/3}$) creates orthogonal triads.
    \item \emph{Real-axis orthogonalities} from $2 \mid n$: the real embedding ($\zeta^{n/2} = -1$) creates orthogonal pairs that interlock the triads.
\end{itemize}

Divisibility by 3 alone (e.g., $n = 9, 15, 21, 27$) creates many triads but they are \emph{completely isolated}---no two share a ray---so the constraint network is satisfiable (the coloring problem decomposes). Divisibility by 2 alone (e.g., $n = 4, 8, 10$) creates pairs but too few interlocking triads. Only when both are present ($6 \mid n$) does the interlocking force a KS contradiction.

\begin{remark}
This result connects to the finding of Cortez, Morales, and Reyes~\cite{CortezMoralesReyes2022} that $\Z[1/6]$ is the minimal ring extension of $\Z$ exhibiting KS contextuality in dimension 3. The prime factorization $6 = 2 \times 3$ appears in both results: the divisibility condition $6 \mid n$ in the cyclotomic setting mirrors the necessity of inverting both 2 and 3 in the ring-theoretic setting. We can make this connection quantitative using their $N(S) = \operatorname{lcm}\{\|v\|^2 : v \in S\}$ invariant, which determines the ring $\Z[1/N(S)]$ in which the projection matrices $P_v = vv^T/\|v\|^2$ live. Computing $N(S)$ for each island's minimal KS set: Eisenstein $N = 6$ ($= 2 \times 3$), Peres $N = 12$ ($= 2^2 \times 3$), $\Z[\sqrt{-2}]$ $N = 12$, Heegner-7 $N = 12$, CK-31 $N = 30$ ($= 2 \times 3 \times 5$). The Eisenstein island achieves exactly the Cortez--Morales--Reyes minimum $N = 6$, confirming it as the algebraic realization of their minimal ring. Every island requires primes 2 and~3; the integer island additionally needs~5 (from the $\|v\|^2 = 5$ rays in CK-31).
\end{remark}

\subsection{Eisenstein integers}

The Eisenstein integers $\Z[\omega]$ (where $\omega = e^{2\pi i/3}$) produce KS sets with a pool minimum of 33 vectors across all coefficient bounds tested (max coefficient 1--3, squared norm cutoff 3--7), certified by OCUS. Increasing the pool size beyond 57 rays does not reduce the minimum below 33.

\subsection{Complex quadratic rings}\label{sec:complex-quadratic}

Extending the real quadratic survey to complex quadratic rings $\Z[\sqrt{-d}]$ with basic alphabets $\{0, \pm 1, \pm\sqrt{-d}\}$, we tested $d = 1, 2, \ldots, 15$. Most produce colorable ray sets. However, $d = 2$ yields a new island:

\begin{table}[ht]
\centering
\caption{Complex quadratic rings $\Z[\sqrt{-d}]$ with alphabet $\{0, \pm 1, \pm\sqrt{-d}\}$.}
\label{tab:complex-quadratic}
\begin{tabular}{lcccccc}
\toprule
Ring & $d$ & Rays & Pairs & Triads & KS? & Min \\
\midrule
$\Z[i]$ & 1 & 25 & 48 & 6 & No & --- \\
$\Z[\sqrt{-2}]$ & 2 & 49 & 120 & 16 & \textbf{Yes} & \textbf{33} \\
$\Z[\sqrt{-3}]$ & 3 & 49 & 114 & 10 & No & --- \\
$\Z[\sqrt{-5}]$ & 5 & 49 & 114 & 10 & No & --- \\
$\Z[\sqrt{-7}]$ & 7 & 49 & 114 & 10 & No & --- \\
\bottomrule
\end{tabular}
\end{table}

The ring $\Z[\sqrt{-2}]$ with alphabet $\{0, \pm 1, \pm\sqrt{-2}\}$ generates 49 rays, 120 orthogonal pairs, and 16 triads---an uncolorable configuration with minimum KS subset of 33 vectors (certified by OCUS within the pool). The cancellation identity is $(\sqrt{-2})^2 = -2$, which in the Hermitian inner product yields $\overline{\sqrt{-2}} \cdot \sqrt{-2} = |\sqrt{-2}|^2 = 2$, producing the same numerical cancellation $1 + 1 = 2$ that drives the Peres island. Indeed, the ray count, pair count, and triad count of $\Z[\sqrt{-2}]$ exactly match those of $\Z[\sqrt{2}]$ (Table~\ref{tab:alphabets}), and the minimum KS subset size is identical (33). The two islands have isomorphic combinatorial structure despite living in different ambient fields ($\R$ vs.\ $\C$).

We also tested the ring of integers of $\Q(\sqrt{-7})$, namely $\Z[(1+\sqrt{-7})/2]$, with the standard alphabet $\{0, \pm 1, \pm\alpha, \pm\bar\alpha\}$ where $\alpha = (1+\sqrt{-7})/2$.  This produced 145 rays (42 triads) that are KS-uncolorable with pool minimum of 43 vectors---a sixth island with a larger minimum than the 33-vector cluster (Section~\ref{sec:heegner}).

\section{Algebraic Island Classification}\label{sec:islands}

\subsection{Cross-product closure as a classification tool}\label{sec:closure}

An important subtlety is that KS behavior depends on the \emph{choice of generating alphabet}, not merely on the ambient number field. For instance, $\Q(\sqrt{5}) = \Q(\varphi)$ as fields, but the alphabet $\{0, \pm 1, \pm\sqrt{5}\}$ produces a colorable set while $\{0, \pm 1, \pm\varphi\}$ (where $\varphi = (1+\sqrt{5})/2$ is the ring-of-integers generator for $\Q(\sqrt{5})$) produces a KS-uncolorable set after completion. We therefore define:

We call the KS-supporting configurations ``islands'' because they are isolated in the space of algebraic parameters: perturbing the generator away from its exact value (e.g., moving $\sqrt{2}$ to $\sqrt{2.01}$, or shifting $\omega$ off the unit circle) immediately destroys KS-uncolorability, and no continuous path through parameter space connects one KS-supporting alphabet to another.  The trigonometric sweep of Remark~\ref{rem:trig} makes this vivid: KS-uncolorability occurs at isolated algebraic angles with colorable seas between them.

\begin{definition}
An \emph{algebraic island} is a pair $(\mathcal{R}, \mathcal{A})$ consisting of a subring $\mathcal{R} \subseteq \C$ (typically the ring of integers of a number field) and a finite alphabet $\mathcal{A} \subset \mathcal{R}$ such that:
\begin{enumerate}[label=(\roman*)]
    \item the ray set $S(\mathcal{A})$ generates vectors in $\mathcal{R}^3$;
    \item the completion of $S(\mathcal{A})$ (see below) terminates in a \emph{finite} ray pool $\bar{S}$ with coordinates in $\mathcal{R}$;
    \item the completed ray pool $\bar{S}$ is KS-uncolorable.
\end{enumerate}
Two islands $(\mathcal{R}_1, \mathcal{A}_1)$ and $(\mathcal{R}_2, \mathcal{A}_2)$ are \emph{equivalent} if their completed ray pools have isomorphic orthogonality hypergraphs.
\end{definition}

The \emph{completion} in condition (ii) depends on whether the island is real or complex:

\begin{itemize}
    \item \textbf{Real islands} ($\mathcal{R} \subset \R$): cross-product completion (Definition~\ref{def:completion}). For each orthogonal pair $v \perp w$, adjoin $[v \times w]$.
    \item \textbf{Complex islands} ($\mathcal{R} \subset \C$, $\mathcal{R} \not\subset \R$): for each Hermitian-orthogonal pair $v \perp w$ in $\C^3$, adjoin the unique ray $[u]$ spanning $(\mathrm{span}\{v,w\})^\perp$. Concretely, $u$ is the vector of $2 \times 2$ cofactors of the matrix $\begin{pmatrix} \bar{v} \\ \bar{w} \end{pmatrix}$:
    \[
    u_k = (-1)^{k+1} \det \begin{pmatrix} \bar{v}_{k+1} & \bar{v}_{k+2} \\ \bar{w}_{k+1} & \bar{w}_{k+2} \end{pmatrix}, \quad k = 1,2,3 \pmod{3}.
    \]
    This is the Hermitian analogue of the cross product. When $\mathcal{R}$ is closed under conjugation, $u \in \mathcal{R}^3$, so completion preserves the coordinate ring.
\end{itemize}

\noindent In $\R^3$, conjugation is trivial and the cofactor formula reduces to the standard cross product, so the two definitions agree.

Note that condition (ii) is automatically satisfied when $\mathcal{R}$ is closed under the coordinate operations of the completion (cross products preserve $\mathcal{R}^3$ for real rings; orthogonal complements preserve $\mathcal{R}^3$ for complex rings closed under conjugation). The \emph{finiteness} of the completed pool and the \emph{combinatorial structure} it generates are the nontrivial properties. The same ring $\mathcal{R}$ may support zero, one, or multiple islands depending on the alphabet.

\begin{table}[ht]
\centering
\caption{Cross-product closure analysis for \emph{real} islands ($\mathcal{R} \subset \R$). ``Compl.''\ = completed ray count; ``New mag.''\ = distinct new coordinate magnitudes introduced by completion (e.g., for integers, cross products introduce 3, 4, 5 beyond $\{0,1,2\}$); ``Min'' = OCUS-certified pool minimum.}
\label{tab:closure}
\small
\begin{tabular}{lcccccc}
\toprule
Alphabet $(\mathcal{R}, \mathcal{A})$ & Raw & Compl. & New mag. & $2{:}1$? & Uncol.? & Min \\
\midrule
$(\Z, \{0,\pm 1,\pm 2\})$ & 49 & 109 & $+3$ & Yes & Yes & \textbf{31} \\
$(\Z[\sqrt{2}], \{0,\pm 1,\pm\!\sqrt{2}\})$ & 49 & 145 & $+12$ & No & Yes & \textbf{33} \\
$(\Z[\sqrt{3}], \{0,\pm 1,\pm\!\sqrt{3}\})$ & 49 & 145 & $+12$ & No & No & --- \\
$(\Z[\sqrt{5}], \{0,\pm 1,\pm\!\sqrt{5}\})$ & 49 & 145 & $+12$ & No & No & --- \\
$(\Z[\varphi], \{0,\pm 1,\pm\varphi\})$ & 49 & 205 & $+21$ & No & Yes & \textbf{52} \\
\bottomrule
\end{tabular}

\medskip
\noindent For the \emph{complex} islands, completion uses the Hermitian orthogonal-complement operation (Section~\ref{sec:closure}). The Eisenstein alphabet $\{0, \pm 1, \pm\omega, \pm\bar\omega\}$ with $\omega = e^{2\pi i/3}$ generates 57 rays (22 triads) at max coefficient 1. Complex completion does not expand this pool (all Hermitian-orthogonal complements already lie within it), and the completed set is KS-uncolorable with pool minimum of 33 vectors (OCUS-certified). The complex quadratic alphabet $\{0, \pm 1, \pm\sqrt{-2}\}$ generates 49 rays (16 triads) with identical combinatorial structure to the Peres island; its pool minimum is also 33 vectors. The Heegner-7 alphabet $\{0, \pm 1, \pm\alpha, \pm\bar\alpha\}$ with $\alpha = (1+\sqrt{-7})/2$ generates 145 rays (42 triads); its pool minimum is 43 vectors (Section~\ref{sec:heegner}).
\end{table}

\subsection{Mixed alphabet evidence}\label{sec:mixed}

Cross-product completion of mixed alphabets---combining rays from two algebraically independent coordinate alphabets, then closing under cross products---provides further evidence for the primacy of the cancellation identity. The cross products generate rays in \emph{neither} original alphabet, exploring configurations invisible to any single-alphabet search.

\begin{table}[ht]
\centering
\caption{Mixed alphabet search with cross-product completion.}
\label{tab:mixed}
\begin{tabular}{lcccc}
\toprule
Combination & Total rays & Triads & KS? & Min \\
\midrule
$\{0,\pm 1,\pm 2\} + \{0,\pm 1,\pm\sqrt{2}\}$ & 205 & 158 & Yes & 31 \\
$\{0,\pm 1,\pm 2\} + \{0,\pm 1,\pm\sqrt{3}\}$ & 217 & 164 & Yes & 31 \\
$\{0,\pm 1,\pm 2\} + \{0,\pm 1,\pm\sqrt{5}\}$ & 217 & 164 & Yes & 31 \\
$\{0,\pm 1,\pm 2\} + \{0,\pm 1,\pm\varphi\}$ & 241 & 188 & Yes & 34$^\dagger$ \\
$\{0,\pm 1,\pm\sqrt{2}\} + \{0,\pm 1,\pm\sqrt{3}\}$ & 229 & 166 & Yes & 33 \\
$\{0,\pm 1,\pm\sqrt{2}\} + \{0,\pm 1,\pm\varphi\}$ & 253 & 190 & Yes & 33 \\
\bottomrule
\end{tabular}
\end{table}

Cross-product completion adds 150+ rays beyond the original alphabets, but the minimum KS size does not improve. The 31-vector barrier persists across all algebraically enriched constructions containing the 2:1 ratio. Conversely, combinations lacking the 2:1 ratio achieve at best 33 despite their larger pools, confirming that the 2:1 ratio is the rate-limiting algebraic feature, not the pool size.

\smallskip\noindent$^\dagger$The pool $\{0,\pm 1,\pm 2\} + \{0,\pm 1,\pm\varphi\}$ contains all 49 integer-alphabet rays as a subset, so it includes CK-31 as a sub-configuration. The reported minimum of 34 reflects the fact that our randomized greedy minimization (200 trials) did not locate the 31-ray subset within the larger 241-ray pool.

\begin{remark}[Relation to the SI-C closure of Trandafir and Cabello]
Trandafir and Cabello~\cite{TrandafirCabello2025} construct their KS sets by starting from the 13-element minimal state-independent contextuality (SI-C) set of Yu and Oh, completing to orthonormal bases (25 rays), and adding all vectors orthogonal to $\geq 2$ existing vectors (97 rays total).  We computed the overlap between this 97-ray SI-C closure and our 49-ray integer pool $\{0,\pm1,\pm2\}$: they share 37 rays, with 60 rays in the SI-C closure only (coordinates up to $\pm 5$, norms up to~35) and 12 rays in the integer pool only (those with two $\pm 2$ coordinates, norm~9).  A second round of orthogonal closure grows the SI-C set to 1{,}741~rays that \emph{contain} the entire integer pool as a subset.  The two constructions---their top-down closure from SI-C and our bottom-up alphabet enumeration---thus converge: CK-31 lives in the intersection of both frameworks.
\end{remark}

\subsection{The six known islands among tested fields}\label{sec:sixislands}

\begin{observation}\label{obs:islands}
Among all alphabets and fields tested in our survey---quadratic fields $\Q(\sqrt{d})$ for $d = 2, \ldots, 30$, cyclotomic fields for $n \leq 30$, and selected extensions---six algebraic islands support KS sets:
\begin{enumerate}[label=(\arabic*)]
    \item \textbf{Integer island} $(\Z, \{0,\pm 1,\pm 2\})$: cancellation $1 + 1 = 2$; pool minimum: 31 vectors ($\CK$). Supported by exhaustive enumeration over the 37-ray union of discovered 31-sets (Proposition~\ref{prop:31optimal}).
    \item \textbf{Peres island} $(\Z[\sqrt{2}], \{0,\pm 1,\pm\sqrt{2}\})$: cancellation $(\sqrt{2})^2 = 1 + 1$; pool minimum: 33 vectors.
    \item \textbf{Eisenstein island} $(\Z[\omega], \{0,\pm 1,\pm\omega,\pm\bar\omega\})$: cancellation $1 + \omega + \omega^2 = 0$; pool minimum: 33 vectors. Identified with Cabello's ``simplest KS set''~\cite{Cabello2025simplest} (Section~\ref{sec:cabello}).
    \item \textbf{Complex quadratic island} $(\Z[\sqrt{-2}], \{0,\pm 1,\pm\sqrt{-2}\})$: cancellation $|\sqrt{-2}|^2 = 2$; pool minimum: 33 vectors. Not a new orthogonality graph: this is a complex coordinatization of the same graph type as the Peres island (Proposition~\ref{prop:isomorphism}), demonstrating that the modulus-2 cancellation mechanism extends to imaginary quadratic rings.
    \item \textbf{Heegner-7 island} $(\Z[(1+\sqrt{-7})/2], \{0,\pm 1,\pm\alpha,\pm\bar\alpha\})$ where $\alpha = (1+\sqrt{-7})/2$: cancellation $\alpha\bar\alpha = 2$; pool minimum: 43 vectors. Newly discovered (Section~\ref{sec:heegner}).
    \item \textbf{Golden island} $(\Z[\varphi], \{0,\pm 1,\pm\varphi\})$: cancellation $\varphi^2 = \varphi + 1$; pool minimum: 52 vectors.
\end{enumerate}
No other alphabet from our survey (real quadratic fields $\Q(\sqrt{d})$ for $d = 2,\ldots,30$; all nine imaginary quadratic fields with class number 1; roots of unity with $6\nmid n$; cubic extensions; icosahedral symmetry directions) produced a KS-uncolorable set.  We emphasize that this is a computational survey over a finite collection of tested families, not a classification of all number fields or all possible two-element alphabets.  The hierarchy $31 < 33 = 33 = 33 < 43 < 52$ gives the exact minimum KS subset size \emph{within each pool}, certified by an OCUS (Optimal Constrained Unsatisfiable Subset) procedure that exhaustively proves no smaller KS subset exists within the respective ray pool (the first five pools in under 4\,s each; the Golden pool in 303\,s).
\end{observation}

This hierarchy correlates with the algebraic complexity of the cancellation identity (Table~\ref{tab:cancellation}).

\subsection{Discovery of the golden ratio island}\label{sec:golden}

The golden ratio field $\Q(\varphi)$ where $\varphi = (1 + \sqrt{5})/2$ presents a remarkable case. The raw alphabet $\{0, \pm 1, \pm\varphi\}$ generates 49 rays with only 10 triads---colorable (the SAT instance is satisfiable). However, cross-product completion generates 156 additional rays (all with coordinates in $\Q(\varphi)$), expanding the pool to 205 rays with 166 triads. This completed set \emph{is} KS-uncolorable, with the smallest KS subset found containing 52 vectors.

Note that $\varphi = (1+\sqrt{5})/2$ is the ring-of-integers generator for $\Q(\sqrt{5})$ (since $5 \equiv 1 \pmod{4}$). The alternative alphabet $\{0, \pm 1, \pm\sqrt{5}\}$ (using the field generator rather than the ring-of-integers generator) produces a \emph{different} and colorable set. This demonstrates that the choice of alphabet within a fixed number field is consequential (see Section~\ref{sec:closure}).

\begin{observation}\label{obs:golden}
The golden ratio island is invisible to standard alphabet searches because the raw alphabet $\{0, \pm 1, \pm\varphi\}$ is colorable. It becomes KS-uncolorable only through the emergent orthogonalities created by cross-product completion, which generates additional rays with coordinates in $\Z[\varphi]$ (the ring of integers of $\Q(\sqrt{5})$). This suggests there may be other algebraic substrates whose KS structure is hidden from raw-alphabet searches and revealed only by completion---though our survey of quadratic and selected higher extensions found no further examples.  We note that golden-ratio coordinates also appear in higher dimensions: Pavicic et al.~\cite{Pavicic2004} catalog KS sets in $d = 4$ with vector components from $\{0, \pm(\sqrt{5}-1)/2, \pm 1, \pm(\sqrt{5}+1)/2, 2\}$, suggesting that $\Q(\varphi)$ is a natural algebraic substrate for KS constructions across multiple dimensions.
\end{observation}

\subsection{Discovery of the Heegner-7 island}\label{sec:heegner}

A systematic search over the imaginary quadratic fields $\Q(\sqrt{-d})$ with class number 1---the \emph{Heegner numbers} $d = 1, 2, 3, 7, 11, 19, 43, 67, 163$---reveals a clean classification. For $d \equiv 3 \pmod{4}$, the ring of integers is $\Z[(1+\sqrt{-d})/2]$, not $\Z[\sqrt{-d}]$, and the generator $\alpha = (1+\sqrt{-d})/2$ satisfies $\alpha\bar\alpha = (1+d)/4$. The results:

\begin{table}[ht]
\centering
\caption{Imaginary quadratic fields with class number 1 (Heegner numbers). ``Gen.\ norm'' $= |\alpha|^2$ where $\alpha$ is the ring-of-integers generator (cf.\ Section~\ref{sec:prelim}). Only $d = 1, 2, 3, 7$ produce KS sets.}
\label{tab:heegner}
\begin{tabular}{cclcccl}
\toprule
$d$ & Gen.\ norm & Ring & Rays & Triads & Uncol.? & Min KS \\
\midrule
1 & 2 & $\Z[i]+(1+i)$ & 127 & 51 & Yes & 33 \\
2 & 2 & $\Z[\sqrt{-2}]$ & 49 & 16 & Yes & 33 \\
3 & 1 & $\Z[\omega]$ & 57 & 22 & Yes & 33 \\
\textbf{7} & \textbf{2} & $\Z[(1+\sqrt{-7})/2]$ & \textbf{145} & \textbf{42} & \textbf{Yes} & \textbf{43} \\
11 & 3 & $\Z[(1+\sqrt{-11})/2]$ & 145 & 30 & No & --- \\
19 & 5 & $\Z[(1+\sqrt{-19})/2]$ & 145 & 30 & No & --- \\
43 & 11 & $\Z[(1+\sqrt{-43})/2]$ & 145 & 30 & No & --- \\
67 & 17 & $\Z[(1+\sqrt{-67})/2]$ & 145 & 30 & No & --- \\
163 & 41 & $\Z[(1+\sqrt{-163})/2]$ & 145 & 30 & No & --- \\
\bottomrule
\end{tabular}
\end{table}

The pattern is sharp: KS-uncolorability requires the alphabet to support a cancellation identity whose terms have squared modulus at most~2. When $|\alpha|^2 = 2$, the identity $\alpha\bar\alpha = 2$ provides the same modulus-2 cancellation as the integer $1+1=2$ and the Peres $(\sqrt{2})^2 = 2$, but through the complex Hermitian inner product. The Eisenstein case ($d = 3$) has $|\omega|^2 = 1$ but compensates with the three-term phase identity $1 + \omega + \omega^2 = 0$. At $|\alpha|^2 \geq 3$ ($d \geq 11$), the squared moduli are too large for dot-product cancellations in $\C^3$.

The Heegner-7 island is structurally distinct from all other islands. To our knowledge, neither the 43-vector Heegner-7 configuration nor the 52-vector golden ratio configuration has been previously reported in the KS literature; existing systematic searches~\cite{Pavicic2004, LiBrightGanesh2024} have focused on the range up to 31 vectors (see Appendix~\ref{app:rays} for explicit descriptions). The Heegner-7 minimal KS set has 43 vectors arranged in 23 bases, with degree distribution $\{4{:}16, 5{:}20, 6{:}4, 8{:}2, 10{:}1\}$ (five distinct degree types and 12 orbit types). Here the \emph{degree sequence} of a KS set is the list of orthogonality degrees of its rays: each ray's degree is the number of other rays orthogonal to it, and the notation $\{k{:}m\}$ means $m$~rays have degree~$k$. The degree sequence is a graph invariant---two KS sets with different degree sequences cannot be isomorphic---and serves as a first-pass fingerprint for distinguishing structurally different configurations. This is richer than both the Eisenstein 33-set (three degree types, 5 orbits) and the Peres/complex-quadratic 33-sets (two degree types, 5--6 orbits), placing Heegner-7 between the 33-cluster and the Golden island in structural complexity.

\begin{observation}\label{obs:gaussian}
The Gaussian integers $\Z[i]$ with the enriched alphabet $\{0, \pm 1, \pm i, \pm(1+i)\}$ are also KS-uncolorable (min = 33), but the minimal set has degree distribution $\{4{:}30, 8{:}3\}$---identical to the Peres and $\Z[\sqrt{-2}]$ islands. This confirms that the classification is governed by the cancellation identity (here $|1+i|^2 = 2$), not the ambient ring: all three are realizations of the ``modulus-2 cancellation'' and produce minimal KS sets sharing all tested graph invariants (Proposition~\ref{prop:isomorphism}).
\end{observation}

\begin{proposition}\label{prop:isomorphism}
\textbf{(Graph and hypergraph isomorphism of modulus-2 islands.)} The Peres $(\Z[\sqrt{2}])$ and complex quadratic $(\Z[\sqrt{-2}])$ orthogonality graphs are isomorphic, verified by the VF2 algorithm~\cite{Cordella2004VF2} (a complete, backtracking graph isomorphism solver that produces explicit vertex bijections).  The isomorphism holds at both the graph level (orthogonal pair structure) and the hypergraph level (triad structure): the VF2-produced bijection maps triads to triads.  Negative controls confirm the method's discriminating power: norm-${\geq}\,3$ fields $\Q(\sqrt{3})$, $\Q(\sqrt{5})$, and $\Q(\sqrt{-5})$ produce 94-basis pools with different graph structure (all colorable, not KS-uncolorable).
\end{proposition}

This reduces the six known algebraic islands to five distinct orthogonality graph types. The Peres, complex quadratic, and Gaussian islands share a single graph type, unified by the modulus-2 cancellation identity $|\alpha|^2 = 2$ operating in different rings ($\Z[\sqrt{2}]$, $\Z[\sqrt{-2}]$,~$\Z[i]$).

\subsection{Dead ends}

\paragraph{Icosahedral symmetry.} The icosahedron has 31 symmetry-axis directions---6 vertex, 15 edge, 10 face---a tantalizing coincidence with $\CK$. However, these directions produce only 5 triads (colorable). After cross-product completion: 91 rays, 65 triads, still colorable. A heuristic explanation: the icosahedral group has rotation orders 2, 3, 5 (angles 180$^\circ$, 120$^\circ$, 72$^\circ$) but \emph{no 90$^\circ$ rotations}, whereas the orthogonal group of $\CK$ is octahedral, the largest finite subgroup of SO(3) containing 90$^\circ$ rotations.

\paragraph{Higher quadratic fields.} All $\Q(\sqrt{d})$ for $d = 3, \ldots, 30$ (excluding $d = 4, 9, 16, 25$) produce $\leq 10$ triads and are colorable, even after cross-product completion (for $d \neq 5$; $d = 5$ generates the golden ratio island as $\Q(\sqrt{5}) = \Q(\varphi)$).

\section{Discussion}\label{sec:discussion}

\subsection{Why six islands among tested fields?}

Among all tested coordinate alphabets, we observe six algebraic islands.  We emphasize that this is a survey finding, not a completeness theorem: untested families (higher-degree extensions, class-number~$> 1$ fields, multi-element alphabets) could in principle harbor additional islands.  The six arise from two fundamentally different cancellation mechanisms:

\begin{enumerate}
    \item[\textbf{(A)}] \textbf{Norm-2 cancellation}: an algebraic identity producing the value~2 from products of ring elements, enabling three-term dot-product cancellations of the form $a_1 b_1 + a_2 b_2 + a_3 b_3 = 0$ with bounded coordinates. Five islands use this mechanism:
    \begin{itemize}
        \item \textbf{Integer} ($1 + 1 = 2$): additive doubling of units.
        \item \textbf{Peres} ($(\sqrt{2})^2 = 2$), \textbf{Complex quadratic} ($|\sqrt{-2}|^2 = 2$), and \textbf{Gaussian} ($|1+i|^2 = 2$): squaring (or modulus-squaring) produces 2. These three share the same minimal combinatorial structure (degree distribution $\{4{:}30, 8{:}3\}$, min 33) despite living in different rings.
        \item \textbf{Heegner-7} ($\alpha\bar\alpha = 2$ where $\alpha = (1+\sqrt{-7})/2$): modulus-2 cancellation through a quadratic integer satisfying $\alpha^2 - \alpha + 2 = 0$. Despite sharing the modulus-2 identity with the Peres family, the richer alphabet (7 elements vs.\ 5) produces a more complex orthogonality graph (42 triads from 145 rays), yielding a larger minimum KS set of 43 vectors.
    \end{itemize}
    \item[\textbf{(B)}] \textbf{Phase cancellation}: a root-of-unity identity $1 + \omega + \omega^2 = 0$ (where $\omega = e^{2\pi i/3}$), in which three complex unit-norm elements sum to zero. This mechanism produces a qualitatively different three-term cancellation: the vanishing sum involves three equal-magnitude terms rather than two small terms balancing one larger term. The \textbf{Eisenstein island} ($\Z[\omega]$, min 33) is the sole representative.
\end{enumerate}

The \textbf{golden ratio island} ($\varphi^2 = \varphi + 1$, min 52) is a genuine boundary case for the dichotomy. The generator $\varphi$ has $|\varphi|^2 \approx 2.618 > 2$ and is not a root of unity, so the raw alphabet satisfies \emph{neither} mechanism---and indeed the raw set is colorable.  Cross-product completion generates the reciprocal $1/\varphi$ (with $|1/\varphi|^2 < 1$), and the completed coordinate set contains cancellations involving both $\varphi$ and $1/\varphi$ that make the completed pool KS-uncolorable.  Whether this should be classified as ``indirect modulus-2'' (since the completed pool contains effective norm-2 cancellations) or as a third mechanism---\emph{completion-induced cancellation}, where the raw alphabet lacks any low-norm identity but cross-product closure generates one---is partly a matter of convention.  We adopt the former interpretation for simplicity but emphasize that the golden ratio is the one case in our survey where the dichotomy is maintained only by appealing to the completed algebra rather than the generating alphabet.  This is also why the two-element alphabet classification (Observation~\ref{obs:two-element}) applies only to raw alphabets before completion: the golden ratio island tests the limits of the two-mechanism framework.

The efficiency of the mechanism correlates inversely with algebraic complexity: the integer identity produces the densest triad network (ratio 0.53), while the golden ratio produces the sparsest (ratio 0.20).  The Peres/$\Z[\sqrt{-2}]$/Gaussian graph isomorphism (Proposition~\ref{prop:isomorphism}) confirms that the classification is by cancellation identity, not ambient field---as first observed by Gould and Aravind~\cite{GouldAravind2010} for the Peres family.

\subsection{Cancellation classification for two-element alphabets}\label{sec:completeness}

Table~\ref{tab:cancellation} summarizes the cancellation identity, alphabet, and a concrete dot-product example for each island.

\begin{table*}[t]
\centering
\caption{The six algebraic islands: alphabets, cancellation identities, and efficiency hierarchy. Rank orders islands by algebraic simplicity of the cancellation; ``Mag.'' is the number of distinct coordinate magnitudes; ``Min'' is the OCUS-certified pool minimum.}
\label{tab:cancellation}
\footnotesize
\setlength{\tabcolsep}{4pt}
\begin{tabular}{@{}cllllrr@{}}
\toprule
Rank & Island & Alphabet $\mathcal{A}$ & Identity & Type & Mag. & Min \\
\midrule
1 & Integer & $\{0,\pm 1,\pm 2\}$ & $1 + 1 = 2$ & Norm-2 & 2 & 31 \\
2 & Peres & $\{0,\pm 1,\pm\!\sqrt{2}\}$ & $(\sqrt{2})^2 = 2$ & Norm-2 & 2 & 33 \\
2$'$ & Eisenstein & $\{0,\pm 1,\pm\omega,\pm\bar\omega\}$ & $1 + \omega + \omega^2 = 0$ & Phase & 1 & 33 \\
2$''$ & $\Z[\sqrt{-2}]$ & $\{0,\pm 1,\pm\!\sqrt{-2}\}$ & $|\sqrt{-2}|^2 = 2$ & Norm-2 & 2 & 33 \\
2$'''$ & Gaussian & $\{0,\pm 1,\pm i,\pm(1{+}i)\}$ & $|1{+}i|^2 = 2$ & Norm-2 & 2 & 33 \\
3 & Heegner-7 & $\{0,\pm 1,\pm\alpha,\pm\bar\alpha\}$ & $\alpha\bar\alpha = 2$ & Norm-2 & 2 & 43 \\
4 & Golden & $\{0,\pm 1,\pm\varphi\} \to$ compl. & $\varphi^2 = \varphi + 1$ & Indirect & 3 & 52 \\
\bottomrule
\end{tabular}
\end{table*}

We can classify all \emph{primitive three-term cancellation identities} available to two-element alphabets.  Consider an arbitrary alphabet $\{0, \pm 1, \pm x\}$ with $x \notin \{0, \pm 1\}$.  The Hermitian product set $\bar{\mathcal{A}} \cdot \mathcal{A}$ consists of $\{0, \pm 1, \pm x, \pm \bar{x}, \pm |x|^2\}$ (with $\bar{x} = x$ in the real case).  A non-trivial orthogonality $\sum_{k=1}^3 \bar{u}_k v_k = 0$ requires three elements of this product set to sum to zero, with at least one involving $x$ (otherwise the identity is trivial).  Table~\ref{tab:two-element} enumerates all such zero-sums.  Note that this classification applies to the \emph{raw} alphabet before completion; cross-product completion can introduce coordinates outside $\mathcal{A}$, potentially enabling additional cancellation identities in the completed pool (as occurs in the golden ratio case).

\begin{table}[h]
\centering
\caption{All primitive three-term zero-sums from the product set $\bar{\mathcal{A}} \cdot \mathcal{A}$ of a two-element alphabet $\{0, \pm 1, \pm x\}$. Each row is an algebraic constraint on~$x$; every solution corresponds to a known island. This classification applies to the raw alphabet before cross-product completion.}
\label{tab:two-element}
\footnotesize
\setlength{\tabcolsep}{3pt}
\begin{tabular}{@{}llll@{}}
\toprule
Zero-sum pattern & Identity on $x$ & Solution & Island \\
\midrule
$1 + 1 - x = 0$ & $x = 2$ & $x = 2$ & Integer \\
$1 + 1 - |x|^2 = 0$ & $|x|^2 = 2$ & $\sqrt{2},\, \sqrt{-2},\, \ldots$ & Peres/$\Z[\!\sqrt{-2}]$/Heeg.-7 \\
$1 + x - |x|^2 = 0$ & $|x|^2 = 1 + x$ & $x = \varphi$ (real) & Golden \\
$1 + x + \bar{x} = 0$ & $x + \bar{x} = -1$; $x^2{+}x{+}1{=}0$ & $x = \omega = e^{2\pi i/3}$ & Eisenstein \\
$1 - x - \bar{x} = 0$ & $x + \bar{x} = 1$ & $d \equiv 1 \pmod{4}$ & (colorable)$^\dagger$ \\
$|x|^2 + |x|^2 - 1 = 0$ & $|x|^2 = \tfrac{1}{2}$ & $x = 1/\sqrt{2}$ & $\equiv$ Peres (rescaled) \\
$x + x - |x|^2 = 0$ & $|x|^2 = 2x$ & $x = 2$ (real) & $\equiv$ Integer \\
\bottomrule
\end{tabular}

\noindent $^\dagger$The zero-sum $1 - x - \bar{x} = 0$ (i.e., $\mathrm{Tr}(x) = 1$) holds for all imaginary quadratic generators with $d \equiv 1 \pmod{4}$ (e.g., $d = -11, -15, -19, \ldots$).  It produces all-nonzero orthogonal pairs and triads, but these triads are structurally decoupled from the one-zero triads and the full ray set remains colorable---no KS set results.  The absence of type-B (mixed) triads is proved for all integer $N \geq 3$ by SMT verification~\cite{deMouraBjorner2008} over all entry patterns; see also Table~\ref{tab:heegner} (all $d \geq 11$ are colorable).
\end{table}

\begin{observation}\label{obs:two-element}
For any two-element alphabet $\{0, \pm 1, \pm x\}$, every primitive three-term cancellation identity in the product set $\bar{\mathcal{A}} \cdot \mathcal{A}$ that leads to KS-uncolorability reduces to one of the six known islands (up to scaling equivalence).  One additional zero-sum pattern exists ($1 - x - \bar{x} = 0$, requiring $\mathrm{Tr}(x) = 1$) but produces only colorable ray sets.
\end{observation}

\emph{Justification.} The product set $\bar{\mathcal{A}} \cdot \mathcal{A}$ for $\mathcal{A} = \{0, \pm 1, \pm x\}$ contains (up to sign) the values $\{0, 1, x, \bar{x}, |x|^2\}$.  A primitive three-term vanishing sum $t_1 + t_2 + t_3 = 0$ with $t_k \in \bar{\mathcal{A}} \cdot \mathcal{A}$ must use three of these values (with signs).  We enumerate all essentially distinct cases by the number of occurrences of each symbol.  The table (Table~\ref{tab:two-element}) lists all solutions: each row is a distinct cancellation pattern, and for each pattern we solve the resulting equation in $x$ and identify the corresponding island (or note colorability).  No additional patterns exist because there are only $\binom{5+2}{3}$ ways to choose three elements (with repetition) from $\{1, x, \bar{x}, |x|^2\}$, most of which yield no solution or a solution equivalent (by $x \mapsto 1/x$, $x \mapsto \bar{x}$, or rescaling) to a listed row.  The $\mathrm{Tr}(x) = 1$ row arises from the sign pattern $(+1, -1, -1)$ on the triple $\{1, x, \bar{x}\}$, which was previously conflated with the $\mathrm{Tr}(x) = -1$ (Eisenstein) pattern under a sign-and-permutation equivalence that is not valid here: the two sums impose different algebraic constraints.  This classification applies to the raw alphabet \emph{before completion}; it does not rule out the possibility that cross-product completion could introduce qualitatively new cancellation mechanisms, though no such case was observed in our survey.

The completeness does \emph{not} extend to larger alphabets.  A three-element alphabet $\{0, \pm 1, \pm x, \pm y\}$ introduces cross-products $\pm xy$, $\pm \bar{x}y$, etc., enabling cancellation identities such as $1 + xy - |x|^2 = 0$ or $x + y + 1 = 0$ that involve two independent generators.  While our survey has tested many such alphabets (Table~\ref{tab:extended} and Section~\ref{sec:mixed}), the space of multi-element alphabets---particularly those arising from higher-degree number fields with multiple independent generators---is not exhaustively covered.  This is the principal remaining gap in the classification.

\begin{remark}[Trigonometric consistency check]\label{rem:trig}
As an independent test of the island classification, we swept the one-parameter family $\mathcal{A}(\theta) = \{0, \pm 1, \pm\cos\theta, \pm\sin\theta\}$ across 164 angles spanning $(0, \pi/2]$. The Pythagorean identity $\cos^2\theta + \sin^2\theta = 1$ serves as the cancellation mechanism. Exactly 5 distinct alphabets (at isolated angles: $\arctan(1/2)$, $\pi/6$, $\arctan(1/\varphi)$, $\arctan(1/\sqrt{2})$, $\pi/4$) are KS-uncolorable; all reduce to known islands via the modulus-2 mechanism or golden ratio completion. No new islands emerge. Moving $1^\circ$ from any uncolorable angle destroys the property, consistent with the algebraic nature of the norm constraint.
\end{remark}

\subsection{Galois-theoretic reformulation}\label{sec:galois}

The six known islands all live in quadratic extensions of $\Q$ (or $\Q$ itself), each with Galois group $\mathrm{Gal}(K/\Q) \cong \Z/2\Z$.  The non-trivial automorphism $\sigma$ sends $\sqrt{d} \mapsto -\sqrt{d}$.  The two cancellation mechanisms admit a clean reformulation in terms of the Galois norm $\mathrm{N}_{K/\Q}(x) = x \cdot \sigma(x)$ and Galois trace $\mathrm{Tr}_{K/\Q}(x) = x + \sigma(x)$ of the ring-of-integers generator~$x$.

\begin{observation}\label{obs:galois-norm-trace}
For the six known KS-supporting alphabets, the cancellation mechanism is determined by two Galois invariants of the generator:
\begin{itemize}
    \item \emph{Modulus-2 cancellation} $\Leftrightarrow$ $|\mathrm{N}_{K/\Q}(x)| = 2$.  Verified: $\mathrm{N}(\sqrt{2}) = -2$, $\mathrm{N}(\sqrt{-2}) = 2$, $\mathrm{N}((1+\sqrt{-7})/2) = 2$.  For imaginary quadratic fields, the algebraic norm equals the Hermitian squared modulus: $\mathrm{N}(x) = x \bar{x} = |x|^2$.
    \item \emph{Phase cancellation} $\Leftrightarrow$ $\mathrm{N}_{K/\Q}(x) = 1$ and $\mathrm{Tr}_{K/\Q}(x) = -1$.  Verified: $\omega$ satisfies $\omega^2 + \omega + 1 = 0$, so $\mathrm{Tr}(\omega) = \omega + \omega^2 = -1$ and $\mathrm{N}(\omega) = \omega \cdot \omega^2 = 1$.  These conditions characterize roots of the cyclotomic polynomial $\Phi_3(x) = x^2 + x + 1$.
\end{itemize}
No other combination of norm and trace values for a quadratic-integer generator has produced a KS-uncolorable set in our survey.
\end{observation}

This reformulation connects the island classification to Gleason's theorem through a second notion of trace.  Gleason's result constrains probability measures on projection operators via the \emph{matrix trace} $\mu(P) = \mathrm{tr}(\rho P)$; the cancellation mechanisms are determined by the \emph{Galois trace} $\mathrm{Tr}_{K/\Q}(x)$.  The logical chain is: Gleason implies KS (trivially, since $v(n) = \langle n | \rho | n \rangle$ is continuous and cannot be $\{0,1\}$-valued on a connected sphere~\cite{RajanVisser2019}); our island classification then asks \emph{which arithmetics} support the finite combinatorial witnesses of this impossibility.  Both traces encode the arithmetic of the primes 2 and~3: the matrix trace requires denominators divisible by 6 for the projection matrices $P_v = vv^T/\|v\|^2$ to live in the coefficient ring (as shown by Cortez et al.~\cite{CortezMoralesReyes2022}), while the Galois trace and norm determine whether the coordinate ring contains elements whose products can sum to zero in exactly the ways that orthogonality requires.

\smallskip\noindent\textbf{Galois action on the KS hypergraph.}  The automorphism $\sigma$ acts coordinate-wise on $K^3$, and since $\sigma$ is a ring homomorphism, it preserves orthogonality: $\langle v | w \rangle = 0 \Rightarrow \langle \sigma(v) | \sigma(w) \rangle = 0$.  Therefore $\sigma$ is an automorphism of both the orthogonality graph and the basis hypergraph.  This gives an additional discrete symmetry beyond the continuous unitary symmetries, and it may partially explain the universality phenomenon (Section~\ref{sec:csw}): realizations of CK-31 in different islands could be related by Galois conjugation composed with a unitary transformation.

\smallskip\noindent\textbf{Connection to Heegner numbers.}  The imaginary quadratic fields $\Q(\sqrt{-d})$ with class number~1 (unique factorization in $\mathcal{O}_K$) correspond to the nine Heegner numbers $d = 1, 2, 3, 7, 11, 19, 43, 67, 163$.  The three KS-supporting imaginary quadratic islands are $d = 2, 3, 7$---all Heegner numbers.  The remaining Heegner numbers fail for transparent arithmetic reasons: $d = 1$ gives $|\mathrm{N}(i)| = 1$ (norm too small; same orthogonality as the trivial alphabet), and $d \geq 11$ gives $|\mathrm{N}(\sqrt{-d})| = d \geq 11$ (norm too large for low-complexity cancellation).  Among imaginary quadratic fields, the KS-supporting ones are precisely those Heegner-number fields whose ring-of-integers generator has algebraic norm~2 ($d = 2, 7$) or is a primitive cube root of unity ($d = 3$).

\smallskip\noindent\textbf{Cyclotomic Galois structure.}  The 6\,$|\,$\,$n$ theorem (Theorem~\ref{thm:6n}) has a natural Galois-theoretic restatement.  The Galois group $\mathrm{Gal}(\Q(\zeta_n)/\Q) \cong (\Z/n\Z)^*$ must have quotients corresponding to both $\Z/2\Z$ (ensuring $-1 \in S_n$, enabling interlocking) and $\Z/3\Z$ (ensuring $\omega \in S_n$, enabling phase cancellation).  This is equivalent to requiring that $(\Z/n\Z)^*$ surjects onto $\Z/6\Z$, which holds if and only if $6 \mid n$.  The conductor of the abelian extension $\Q(\zeta_6)/\Q$ is~6---the same number appearing in the Cortez--Morales--Reyes ring condition $6 \mid N$.

\smallskip\noindent\textbf{Open directions.}  Several questions arise from this reformulation: (i)~Can the norm/trace characterization of Observation~\ref{obs:galois-norm-trace} be proven exhaustive for quadratic fields, establishing that $|\mathrm{N}| = 2$ and $(\mathrm{N} = 1, \mathrm{Tr} = -1)$ are the \emph{only} norm/trace pairs yielding low-complexity cancellation? (ii)~Does the compositum $\Q(\sqrt{2}, \sqrt{-3})$, with Galois group $(\Z/2\Z)^2$, produce denser KS sets or smaller minimums than either component island? (iii)~For the cubic island $\Q(\sqrt[3]{2})$, does the Galois closure (with group $S_3$) impose additional structural constraints? (iv)~Can class field theory, specifically the conductor of abelian extensions, provide a unified invariant classifying KS-supporting fields?

\subsection{Connection to perfect quantum strategies}\label{sec:cabello}

Our algebraic island framework gains new significance from three recent results of Cabello and collaborators that establish KS sets as the fundamental engine of \emph{perfect quantum strategies}.

Cabello~\cite{Cabello2024bipartite} proved that every bipartite perfect quantum strategy---a strategy achieving the maximum algebraic violation of a Bell inequality in the finite-dimensional projective measurement framework---defines a KS set. More precisely, the projective measurements achieving a perfect strategy on a Bell scenario must contain a KS-uncolorable subset. This means KS sets are not merely foundational curiosities but are \emph{operationally necessary} for the strongest form of quantum advantage in nonlocal games.

The connection deepens: Liu et al.~\cite{Cabello2024equivalences} established a chain of equivalences showing that face nonsignaling correlations, full Bell nonlocality, all-versus-nothing proofs, and pseudotelepathy are all manifestations of the same underlying structure---which, by the first result, is always a KS set. The algebraic islands we classify are therefore not just the algebraic substrates of KS sets; they are the algebraic substrates of \emph{all} perfect quantum strategies in dimension 3.

Most directly relevant to our framework, Cabello~\cite{Cabello2025simplest} constructed what he calls the ``simplest KS set'': a 33-vector, 14-basis configuration in $\C^3$ built from the Eisenstein integers $\Z[\omega]$ via the Weyl--Heisenberg group. This construction has automorphism group of order 144 and is built from the eigenbases of the Weyl--Heisenberg operators $X$ and $Z$ acting on the minimal SI-C set. Crucially, the coordinate entries are drawn from $\{0, 1, \omega, \bar\omega\}$---precisely the Eisenstein alphabet $\{0, \pm 1, \pm\omega, \pm\bar\omega\}$ (after accounting for ray equivalence under scalar multiplication). Cabello's simplest KS set is the smallest member of our Eisenstein island.

This identification provides a concrete bridge between the algebraic island classification and the operational significance of KS sets.

\subsubsection{From KS sets to Bell scenarios: the role of context count}\label{sec:ks-bell}

The operational link between KS sets and nonlocal games is made precise by Trandafir and Cabello~\cite{TrandafirCabello2025optimal}, who give an algorithm to convert any KS set into a bipartite perfect quantum strategy (BPQS) with the minimum number of measurement settings. In such a conversion, the \emph{number of bases} (contexts) in the KS set directly determines the number of inputs for each party in the resulting Bell scenario: fewer bases yield a simpler nonlocal game with fewer measurement choices.

For the Eisenstein island, Cabello~\cite{Cabello2025simplest} shows that the 33-vector, 14-basis KS set yields a $5 \times 9$ qutrit--qutrit BPQS (Alice has 5 inputs, Bob has 9), improving on the $7 \times 9$ strategy derived from the Peres 33-vector set (16 bases).

We independently computed B-KS input counts for all six islands using a SAT-based encoding. Following Definition~2 of Cabello~\cite{Cabello2024bipartite}, a pair $(S_A, S_B)$ of basis subsets is \emph{bipartite KS-uncolorable} (B-KS) if there is no assignment $f: (\mathcal{V}_A \cup \mathcal{V}_B) \times (S_A \cup S_B) \to \{0,1\}$ satisfying:
\begin{itemize}
    \item[(I\,$'$)] \emph{Cross-party exclusion}: $f(u,b) + f(u',b') \leq 1$ for each $b \in S_A$, $b' \in S_B$, $u \in b$, $u' \in b'$ with $u \perp u'$;
    \item[(II\,$'$)] \emph{Completeness}: $\sum_{u \in b} f(u,b) = 1$ for each $b \in S_A \cup S_B$,
\end{itemize}
where $\mathcal{V}_A$ ($\mathcal{V}_B$) is the set of all vectors appearing in $S_A$ ($S_B$). Note that $f$ is contextual: the same vector may receive different values in different bases. Crucially, orthogonality constraints are imposed \emph{only between} $S_A$ and $S_B$ bases, not within a single party's basis set. The basis sets $S_A$ and $S_B$ are allowed to overlap. The optimal BPQS minimizes the product $|S_A| \times |S_B|$ (the total number of input pairs in the Bell scenario). We encode conditions (I$'$) and (II$'$) as a SAT instance and search for the $(S_A, S_B)$ pair minimizing this product, following the approach of Trandafir and Cabello~\cite{TrandafirCabello2025optimal}. For islands with $|\mathcal{B}| \leq 16$ (Eisenstein, Peres, $\Z[\sqrt{-2}]$), we perform exhaustive search over all basis subset pairs; for larger islands (CK-31, Heegner-7, Golden), we use randomized greedy search (500 trials), yielding upper bounds on the optimal product.

\begin{table}[ht]
\centering
\caption{B-KS input counts for all six islands. $|S_A| \times |S_B|$ is the minimum product over bipartite KS-uncolorable basis subset pairs. Eisenstein, Peres, and CK-31 independently verify Cabello~\cite{Cabello2025simplest} Table~I; the last three are new. ``Exact'': exhaustive search; ``best found'': greedy search (500 trials), yielding upper bounds.}
\label{tab:bpqs}
\begin{tabular}{lcccccl}
\toprule
Island & rays & $|\mathcal{B}|$ & $|S_A|$ & $|S_B|$ & Product & Status \\
\midrule
Eisenstein & 33 & 14 & 5 & 9 & \textbf{45} & Exact; verified \\
Peres & 33 & 16 & 7 & 9 & 63 & Exact; verified \\
$\Z[\sqrt{-2}]$ & 33 & 16 & 7 & 9 & 63 & Exact; new ($=$ Peres) \\
Integer (CK-31) & 31 & 17 & 8 & 9 & $\leq$72 & Best found; matches~\cite{Cabello2025simplest} \\
Heegner-7 & 43 & 23 & 9 & 12 & $\leq$108 & Best found; new \\
Golden & 52 & 25 & 12 & 13 & $\leq$156 & Best found; new \\
\bottomrule
\end{tabular}
\end{table}

\noindent Two empirical patterns emerge from Table~\ref{tab:bpqs}. First, $|S_A| + |S_B| \geq |\mathcal{B}|$ in every case examined: the best-found Alice and Bob basis sets nearly \emph{partition} the full basis set (equality holds for Eisenstein, Peres, CK-31, and Golden; Heegner-7 achieves $21/23$). This pattern suggests---but does not prove---that B-KS uncolorability requires the two parties to collectively engage essentially all measurement contexts. Second, the BPQS cost $|S_A| \times |S_B|$ scales roughly as $|\mathcal{B}|^2/4$, with the ratio product$/|\mathcal{B}|^2$ lying between 0.20 (Heegner-7) and 0.25 (Golden) for all six islands. That the Eisenstein island (product$/|\mathcal{B}|^2 = 0.23$, exact) is more efficient than the $\approx 0.25$ achieved by the real-field islands (Peres, CK-31, Golden) may reflect the richer symmetry of $\Z[\omega]$.

The island classification thus gives direct operational meaning to the algebraic substrate. CK-31 achieves the minimum ray count (31) but at the cost of having the most bases among the 33-class islands (17 vs.\ 14). The Eisenstein island sacrifices two rays to gain a sparser basis structure, yielding the simplest Bell scenario. This trade-off---\emph{ray economy vs.\ context economy}---is invisible without the island framework.

The trade-off extends further when CSW invariants (Section~\ref{sec:csw}) and critical bases analysis (Section~\ref{sec:critical-bases}) are included. No single island dominates on all axes: CK-31 is the tightest proof ($\eta = 1.000$) with a nontrivial CSW sandwich, but its 17 bases make it operationally costly. The Eisenstein island has the simplest BPQS and near-maximal tightness ($\eta = 0.929$), but only a marginal CSW advantage. The Heegner-7 island has the richest CSW structure but the largest BPQS cost. This suggests that the choice of algebraic substrate has genuine consequences for quantum information tasks, and that no single notion of ``simplest'' or ``best'' KS set is appropriate---it depends on the operational criterion.

The fact that Cabello's ``simplest'' construction lands precisely on one of our independently discovered islands---and not between them---is strong evidence that the island classification captures genuine algebraic structure rather than an artifact of our search methodology.

\subsection{Graph-theoretic invariants and Bell inequalities}\label{sec:csw}

If the algebraic substrate controls KS-uncolorability, it should also control the \emph{strength} of the contextual advantage. The Cabello--Severini--Winter (CSW) framework~\cite{CabelloSeveriniWinter2014} provides a quantitative test, mapping the orthogonality graph $G$ of a set of rays to a noncontextuality inequality via three invariants:
\[
\alpha(G) \leq \vartheta(G) \leq \alpha^*(G),
\]
where $\alpha(G)$ is the independence number (classical bound), $\vartheta(G)$ is the Lov\'asz theta number (quantum bound, computed by semidefinite programming), and $\alpha^*(G)$ is the fractional packing number (nonsignaling bound, computed by linear programming). A quantum advantage exists when $\vartheta(G) > \alpha(G)$.

We compute these invariants for two representations of each island: the SAT-minimized KS set (minimum cardinality) and the full ray pool (all rays generated by the alphabet before minimization). This comparison reveals a structural mechanism by which cardinality minimization destroys CSW contextual~advantage.

\begin{table}[ht]
\centering
\caption{CSW graph invariants for minimized KS sets. $\alpha$: independence number (classical), $\vartheta$: Lov\'asz theta (quantum, SDP), $\alpha^*$: fractional packing (nonsignaling, LP). Ratio $\vartheta/\alpha$ measures contextual advantage.}
\label{tab:csw-min}
\begin{tabular}{lccccc}
\toprule
Island (minimized) & $n$ & $\alpha(G)$ & $\vartheta(G)$ & $\alpha^*(G)$ & $\vartheta/\alpha$ \\
\midrule
Integer (CK-31) & 31 & 11 & 11.71 & 12.00 & \textbf{1.065} \\
Eisenstein & 33 & 12 & 12.05 & 13.00 & 1.004 \\
Peres & 33 & 12 & 12.00 & 12.00 & 1.000 \\
$\Z[\sqrt{-2}]$ & 33 & 12 & 12.00 & 12.00 & 1.000 \\
Heegner-7 & 43 & 16 & 16.00 & 16.00 & 1.000 \\
\bottomrule
\end{tabular}
\end{table}

\begin{table}[ht]
\centering
\caption{CSW graph invariants for full (non-minimized) ray pools. $\vartheta$ computed via SCS (tolerance $10^{-9}$); $\alpha^*$ via HiGHS (exact). The ``auxiliary'' column counts rays that participate in pairwise orthogonalities but not in any complete basis (Definition~\ref{def:auxiliary}).}
\label{tab:csw-pool}
\begin{tabular}{lccccccc}
\toprule
Island (full pool) & $n$ & bases & auxiliary & $\alpha(G)$ & $\vartheta(G)$ & $\alpha^*(G)$ & $\vartheta/\alpha$ \\
\midrule
Integer & 49 & 26 & 0 (0\%) & 17 & 17.70 & 19.00 & 1.041 \\
Eisenstein & 57 & 22 & 0 (0\%) & 18 & 19.34 & 21.00 & 1.074 \\
Heegner-7 & 145 & 42 & 76 (52\%) & 50 & 55.89 & 66.00 & \textbf{1.118} \\
Peres & 49 & 16 & 16 (33\%) & 23 & 23.00 & 23.00 & 1.000 \\
$\Z[\sqrt{-2}]$ & 49 & 16 & 16 (33\%) & 23 & 23.00 & 23.00 & 1.000 \\
\bottomrule
\end{tabular}
\end{table}

\subsubsection{Auxiliary rays and the CSW advantage}

The comparison between Tables~\ref{tab:csw-min} and~\ref{tab:csw-pool} reveals a striking pattern: the Heegner-7 island's full pool has the strongest contextual advantage of any island ($\vartheta/\alpha = 1.118$), yet its minimized 43-vector set is $\vartheta$-perfect ($\vartheta = \alpha$). To explain this, we introduce a structural distinction among rays in an algebraic pool.

\begin{definition}\label{def:auxiliary}
Let $\mathcal{P}$ be a ray pool with orthogonality graph $G = (V, E)$ and set of bases (maximal orthogonal triples) $\mathcal{B}$. A ray $v \in V$ is \textbf{basis-participating} if $v$ belongs to at least one basis $B \in \mathcal{B}$, and \textbf{auxiliary} otherwise. That is, an auxiliary ray has at least one orthogonal neighbor in $G$ but does not belong to any complete orthonormal triple.
\end{definition}

Auxiliary rays contribute edges to the orthogonality graph without contributing to the hypergraph structure that determines KS-uncolorability. Since KS-uncolorability is a property of the basis hypergraph (no consistent $\{0,1\}$ coloring assigning exactly one ``1'' per basis), auxiliary rays are invisible to it. But they are fully visible to the CSW framework, which depends only on the orthogonality \emph{graph}.

\begin{observation}\label{obs:auxiliary}
\textbf{(Auxiliary rays amplify the CSW advantage.)} Among the five algebraic islands, the CSW contextual advantage $\vartheta(G)/\alpha(G)$ of the full ray pool correlates with the proportion of auxiliary rays:
\begin{center}
\begin{tabular}{lccl}
\toprule
Island & Auxiliary fraction & $\vartheta/\alpha$ (pool) & $\vartheta/\alpha$ (minimized) \\
\midrule
Heegner-7 & 52\% & \textbf{1.118} & 1.000 \\
Peres & 33\% & 1.000 & 1.000 \\
$\Z[\sqrt{-2}]$ & 33\% & 1.000 & 1.000 \\
Integer & 0\% & 1.041 & 1.065 \\
Eisenstein & 0\% & 1.074 & 1.004 \\
\bottomrule
\end{tabular}
\end{center}
Auxiliary rays constrain the independence number $\alpha(G)$ (by adding edges that prevent large independent sets) without contributing to the basis structure. The SDP relaxation $\vartheta(G)$ can exploit these additional constraints by distributing weight fractionally across the auxiliary vertices, widening the gap $\vartheta - \alpha$.
\end{observation}

The mechanism operates as follows. In the Heegner-7 pool ($n = 145$), only 69 rays participate in the 42 bases. The remaining 76 auxiliary rays add 522 orthogonal pairs to the graph, including 9 hub vertices of degree 16 (each connected to 11\% of all other vertices). These hubs severely restrict the size of independent sets: only $\alpha/n = 0.345$ of vertices can be simultaneously independent, compared to $0.469$ for the Peres and $\Z[\sqrt{-2}]$ pools (which are $\vartheta$-perfect). The SDP, however, can assign fractional weights summing to $\vartheta/n = 0.386$, yielding a 12\% quantum advantage.

When SAT minimization reduces the Heegner-7 pool to 43 vectors, it removes auxiliary rays first---they are dispensable for KS-uncolorability. But these are precisely the rays that constrained $\alpha$ and enabled the $\vartheta > \alpha$ gap. The minimized set is $\vartheta$-perfect.

\subsubsection{Further patterns}

\begin{enumerate}
    \item \textbf{Full pools consistently outperform minimized sets} for the CSW advantage. The Eisenstein pool jumps from $\vartheta/\alpha = 1.004$ (minimized) to $1.074$ (full pool). The integer pool is the exception: it drops slightly from $1.065$ to $1.041$, suggesting that CK-31 is unusually efficient at preserving CSW structure despite minimization. Notably, the integer and Eisenstein pools have zero auxiliary rays---every ray participates in a basis---so their CSW advantages arise from graph regularity rather than auxiliary edges.

    \item \textbf{The nonsignaling bound is tightly constrained by basis structure.} For the four minimized sets where $\vartheta = \alpha$ (Peres, $\Z[\sqrt{-2}]$, Heegner-7, plus Eisenstein near-equality), we also have $\alpha^* = \alpha$: all three invariants collapse. Only the integer island (CK-31) maintains a nontrivial sandwich $\alpha = 11 < \vartheta = 11.71 < \alpha^* = 12$. In the full pools, the Heegner-7 pool shows the widest separation: $\alpha = 50 < \vartheta = 55.89 < \alpha^* = 66$, with the nonsignaling bound 32\% above the classical bound.

    \item \textbf{KS-uncolorability does not imply CSW violation}: The Peres and $\Z[\sqrt{-2}]$ islands remain $\vartheta$-perfect even in their full pools, despite having 16 auxiliary rays each. The auxiliary ray mechanism is necessary but not sufficient; the specific placement of orthogonalities matters. In these two pools, the orthogonality graph is disconnected: the 16 auxiliary rays (degree 3) form loosely connected components that do not bridge the main structure. For disconnected graphs, $\vartheta(G) = \sum_i \vartheta(G_i)$ and $\alpha(G) = \sum_i \alpha(G_i)$, so $\vartheta$-perfection of each component implies $\vartheta$-perfection of the whole~\cite{Knuth1994}. This contrasts with the Heegner-7 and integer pools, where auxiliary rays create a connected graph with enough edge density to separate $\vartheta$ from $\alpha$.
\end{enumerate}

\begin{remark}
The distinction between basis-participating and auxiliary rays reveals that KS-uncolorability and CSW contextual advantage are driven by different structural features of the same algebraic pool. Uncolorability requires bases arranged so that no valid $\{0,1\}$ coloring exists---a property of the basis \emph{hypergraph}. The CSW advantage requires pairwise orthogonalities arranged so that $\vartheta > \alpha$---a property of the orthogonality \emph{graph}. Auxiliary rays contribute to the latter without affecting the former. The operational lesson is that the contextual strength of an algebraic island, in the CSW sense, is determined by its full ray pool, not its minimal KS subset. Minimizing for cardinality can destroy the CSW advantage entirely.
\end{remark}

\subsubsection{Spectral explanation of the CSW advantage}\label{sec:spectral}

The Hoffman bound provides a spectral explanation for which pools admit a CSW violation. For any graph $G$ with adjacency matrix $A$ having largest eigenvalue $\lambda_{\max}$ and smallest eigenvalue $\lambda_{\min}$, the independence number satisfies~\cite{Haemers1995}
\[
\alpha(G) \leq \frac{n \cdot (-\lambda_{\min})}{\lambda_{\max} - \lambda_{\min}}.
\]
For connected graphs, this bound is tight and well-established~\cite{Haemers1995}. However, for disconnected graphs $\alpha(G)$ can exceed $H(G)$, as Table~\ref{tab:spectral} confirms for the Peres and $\Z[\sqrt{-2}]$ pools ($\alpha/H = 1.176$), since the global eigenvalue extrema may not reflect the per-component structure. When $\alpha(G)$ is well below this bound, the graph is ``spectrally tight'': the eigenvalue structure constrains independent sets more than the combinatorial structure alone requires, leaving room for the SDP relaxation $\vartheta(G)$ to exceed $\alpha(G)$.

\begin{table}[ht]
\centering
\caption{Spectral properties of full ray pool orthogonality graphs. $H$ = Hoffman bound on $\alpha$. The ratio $\alpha/H$ measures spectral tightness: values near or below 1 correlate with CSW violations.}
\label{tab:spectral}
\begin{tabular}{lccccccc}
\toprule
Pool & $n$ & $\lambda_{\max}$ & $\lambda_{\min}$ & $H$ & $\alpha$ & $\alpha/H$ & $\vartheta/\alpha$ \\
\midrule
Heegner-7 & 145 & 8.25 & $-4.50$ & 51.2 & 50 & \textbf{0.977} & \textbf{1.118} \\
Eisenstein & 57 & 6.42 & $-3.45$ & 19.9 & 18 & \textbf{0.904} & \textbf{1.074} \\
Integer & 49 & 5.98 & $-3.29$ & 17.4 & 17 & 0.977 & 1.041 \\
Peres & 49 & 5.53 & $-3.68$ & 19.6 & 23 & 1.176 & 1.000 \\
$\Z[\sqrt{-2}]$ & 49 & 5.53 & $-3.68$ & 19.6 & 23 & 1.176 & 1.000 \\
\bottomrule
\end{tabular}
\end{table}

The pattern in Table~\ref{tab:spectral} is clear: the three pools with $\alpha/H < 1$ (spectrally tight) have CSW violations, with the largest violation ($\vartheta/\alpha = 1.118$) occurring for the Heegner-7 pool ($\alpha/H = 0.977$). The two $\vartheta$-perfect pools have $\alpha/H > 1$, meaning $\alpha$ exceeds the Hoffman bound; this is expected rather than anomalous, because these graphs are disconnected (the 16 auxiliary rays of degree 3 form loosely connected components) and the Hoffman bound does not apply to disconnected graphs in its standard form.

\subsection{Contextual tightness: critical bases analysis}\label{sec:critical-bases}

The CSW framework measures how much quantum mechanics exceeds classical bounds on the orthogonality \emph{graph}. We now introduce a complementary measure that probes the internal structure of the basis \emph{hypergraph}: how many bases must be removed before a KS set becomes colorable?

\begin{definition}\label{def:essential-basis}
Let $\mathcal{S}$ be a KS set with basis set $\mathcal{B} = \{B_1, \ldots, B_m\}$. A basis $B_i$ is \textbf{essential} if $\mathcal{S}$ with the constraint from $B_i$ removed admits a valid $\{0,1\}$ coloring---i.e., the set is KS-uncolorable only when $B_i$'s ``exactly one green'' rule is enforced. The \textbf{critical number} $\kappa(\mathcal{S})$ is the minimum $k$ such that some $k$-element subset of bases can be removed to make $\mathcal{S}$ colorable. The \textbf{contextual fraction} is $\kappa(\mathcal{S}) / |\mathcal{B}|$.
\end{definition}

Table~\ref{tab:critical-bases} reports these quantities for each island's SAT-minimized KS set and the full Heegner-7 pool, where $\eta = (\text{essential bases}) / |\mathcal{B}|$.

\begin{table}[ht]
\centering
\caption{Critical bases analysis. $|\mathcal{B}|$ = number of bases, ess = essential bases (whose individual removal breaks uncolorability), $\eta$ = essentiality ratio, $\kappa$ = critical number (minimum bases to remove), CF = contextual fraction $\kappa/|\mathcal{B}|$.}
\label{tab:critical-bases}
\begin{tabular}{lcccccc}
\toprule
Set & rays & $|\mathcal{B}|$ & ess & $\eta$ & $\kappa$ & CF \\
\midrule
Integer (CK-31) & 31 & 17 & \textbf{17} & \textbf{1.000} & 1 & 0.059 \\
Eisenstein min-33 & 33 & 14 & 13 & 0.929 & 1 & 0.071 \\
Peres min-33 & 33 & 16 & 13 & 0.812 & 1 & 0.063 \\
$\Z[\sqrt{-2}]$ min-33 & 33 & 16 & 13 & 0.812 & 1 & 0.063 \\
Heegner-7 min-43 & 43 & 23 & 18 & 0.783 & 1 & 0.044 \\
\midrule
Peres pool & 49 & 16 & 13 & 0.812 & 1 & 0.063 \\
Heegner-7 pool & 145 & 42 & 0 & 0.000 & \textbf{2} & 0.048 \\
\bottomrule
\end{tabular}
\end{table}

CK-31 is maximally tight ($\eta = 1.000$): every basis is essential.  Peres and $\Z[\sqrt{-2}]$ are again identical (16 bases, 13 essential, $\eta = 0.812$), extending graph isomorphism to the hypergraph level.  The Heegner-7 full pool is uniquely robust ($\kappa = 2$): no single basis removal breaks uncolorability, but 24 of 861 basis pairs are critical.  The transition from $\kappa = 1$ (all minimized sets) to $\kappa = 2$ shows that algebraic completeness provides contextual redundancy that minimization destroys.

\subsection{Supporting evidence}\label{sec:supporting}

\begin{definition}\label{def:merge}
Let $G = (V, E, T)$ be a KS geometry with vertex set $V$ (rays), edge set $E$ (orthogonal pairs), and triad set $T$ (orthogonal triples).  A \emph{merge} of a non-orthogonal pair $\{v_1, v_2\} \notin E$ is the geometry $G' = (V', E', T')$ obtained by: (1)~removing $v_2$ from $V$; (2)~deleting all edges incident to $v_2$ from $E$; (3)~reconstructing triads $T'$ as all 3-cliques in the resulting edge set $E'$.  The merge is \emph{KS-preserving} if $G'$ is KS-uncolorable.  A KS set is \emph{merge-saturated} if every non-orthogonal pair merge is KS-preserving.
\end{definition}

\begin{observation}[Universal merge saturation]\label{obs:merge}
For each of the six minimal KS sets, we tested every non-orthogonal vertex pair merge (Definition~\ref{def:merge}). In all 3{,}756 merges across the six islands---394 (CK-31), 456 (Peres-33), 450 (Eisenstein-33), 456 ($\Z[\sqrt{-2}]$-33), 796 (Heegner-7-43), 1{,}204 (Golden-52)---the merged $(n{-}1)$-vertex graph remains KS-uncolorable, verified by SAT. We conjecture that every minimal KS set in $\R^3$ or $\C^3$ is merge-saturated.
\end{observation}

\begin{observation}[Rigidity of the six islands]\label{obs:rigidity}
Trandafir and Cabello~\cite{TrandafirCabello2025} define a KS set as \emph{rigid} if its orthogonality graph admits a unique geometric realization up to unitary equivalence in~$\C^3$: any set of rank-1 projectors satisfying the same graph relations must be unitarily equivalent to the original.  They proved that $\CK$ is rigid and that Peres-33 is not rigid (the Penrose~33-set provides an inequivalent realization of the same graph).

We test rigidity for all six minimal KS sets by computing the Jacobian of the constraint system (normalization $\|v_i\|^2 = 1$ and orthogonality $\braket{v_i}{v_j} = 0$ for all edges) at the known solution and comparing the null-space dimension with the expected symmetry dimension.  For $n$ unit vectors in~$\C^3$, the symmetry group---$U(3)$ acting globally plus $U(1)^n$ acting ray-by-ray, modulo the overlap of overall phase---has dimension $n + 8$.  A set is \emph{infinitesimally rigid} if the null space of the Jacobian has dimension exactly $n + 8$ (no deformation directions beyond symmetries).  For $\CK$, Trandafir and Cabello~\cite{TrandafirCabello2025} proved \emph{global} rigidity (unique up to unitary equivalence); for the remaining islands, our Jacobian test establishes infinitesimal rigidity, which is necessary but not sufficient for global rigidity.

\medskip
\begin{center}
\renewcommand{\arraystretch}{1.1}
\begin{tabular}{lccccc}
\hline
Island & $n$ & Pairs & $\dim\ker J$ & $n{+}8$ & Status \\
\hline
CK-31 (integer)   & 31 & 71  & 39 & 39 & inf.\ rigid \\
Eisenstein-33      & 33 & 78  & 41 & 41 & inf.\ rigid \\
Peres-33           & 33 & 72  & 42 & 41 & flex (1) \\
$\Z[\sqrt{-2}]$-33 & 33 & 72 & 42 & 41 & flex (1) \\
Heegner-7          & 43 & 107 & 51 & 51 & inf.\ rigid \\
Golden-52          & 52 & 124 & 60 & 60 & inf.\ rigid \\
\hline
\end{tabular}
\end{center}
\medskip

Four of the six islands are infinitesimally rigid in~$\C^3$; the Peres and $\Z[\sqrt{-2}]$ islands each have exactly one deformation dimension.  For $\CK$, infinitesimal rigidity is strengthened to global rigidity by Trandafir and Cabello~\cite{TrandafirCabello2025}.  For Eisenstein-33, Cabello~\cite{Cabello2025simplest} independently claims rigidity (unique up to unitary transformation, proved using the method of~\cite{TrandafirCabello2025}); our Jacobian analysis provides concordant evidence via an independent method.  For Eisenstein-33, Heegner-7, and Golden-52, we searched for distinct geometric realizations by random perturbation of the known solution followed by constrained optimization back to the orthogonality surface (20 trials per island, perturbation scale~$0.3$); in every trial the optimizer converged to a realization unitarily equivalent to the original, providing computational evidence consistent with global rigidity.  A formal proof (e.g., via Gr\"obner bases as in~\cite{TrandafirCabello2025}) remains open.

The structural explanation is clean: the Peres and $\Z[\sqrt{-2}]$ islands share the same orthogonality graph (Proposition~\ref{prop:isomorphism}), and the flex belongs to the \emph{graph}, not the algebra---embedding the Peres vectors (all real) into~$\C^3$ produces the same one-dimensional flex.  The Eisenstein graph has 78 orthogonal pairs (vs.\ 72 for Peres), including 12 vertices of degree~5 (vs.\ all degree~4 in Peres); those 6~extra pairs provide 12~additional constraints that kill the flex.

The Peres/\mbox{$\Z[\sqrt{-2}]$} flex is \emph{infinitesimal} only: second-order obstruction analysis shows a nonzero component in the cokernel of the Jacobian, blocking extension to a finite deformation.  The set is infinitesimally flexible but finitely rigid---the moduli space has isolated points (Peres and Penrose realizations) rather than a continuous family.
\end{observation}

\begin{remark}[Algebraic alphabets as a new construction family]\label{rem:new-construction}
Trandafir and Cabello~\cite{TrandafirCabello2025} identify two construction families for KS sets in~$\C^3$---dimension-lifting and concatenation---and show that neither produces rigid sets.  The algebraic alphabet method falls outside both families: rays are pinned to algebraic values by number-theoretic constraints rather than assembled from flexible geometric pieces.  Indeed, four of six islands are rigid in~$\C^3$ (Observation~\ref{obs:rigidity}), showing that the method extends the known construction landscape beyond the two families analyzed in~\cite{TrandafirCabello2025}.
\end{remark}

\begin{observation}[Three distinct 33-vector KS sets]\label{obs:three-33}
The literature now contains three minimal 33-vector KS sets with pairwise non-isomorphic orthogonality graphs:

\medskip
\begin{center}
\renewcommand{\arraystretch}{1.1}
\begin{tabular}{lcccccc}
\hline
Set & Field & Pairs & Triads & Degree range & Inf.\ rigid \\
\hline
CK-33 (Conway--Kochen) & $\Z$ & 76 & 20 & $3$--$8$ & yes \\
Eisenstein-33 (this paper) & $\Z[\omega]$ & 78 & 14 & $4$--$8$ & yes \\
Peres-33 & $\Z[\sqrt{2}]$ & 72 & 16 & $4$--$8$ & no \\
\hline
\end{tabular}
\end{center}
\medskip

Graph non-isomorphism is confirmed by VF2 and by the distinct degree sequences: CK-33 has eight vertices of degree~3 and two of degree~8; Eisenstein-33 has twelve vertices of degree~5 and three of degree~8; Peres-33 has thirty vertices of degree~4 and three of degree~8.  No pair of these three graphs is isomorphic.

CK-33 was communicated by Conway and Kochen to Peres circa 1990~\cite{Peres1991}; its 33~vectors use the same integer alphabet $\{0, \pm 1, \pm 2\}$ as $\CK$ and lie entirely within the 49-ray integer pool.  However, CK-33 is \emph{minimal at~33 within its own subgraph}: greedy deletion cannot reduce it below 33~vectors, confirming that it is a genuinely distinct minimal KS set, not a non-minimal subset of $\CK$.  The integer pool thus contains at least two structurally different KS sets---one at 31~vectors (the $\CK$ family of six MUS sets; Observation~\ref{obs:mus-landscape}) and one at~33 (CK-33).

The three 33-sets also exhibit a striking triad-vs-pair tradeoff: CK-33 has the most triads (20) but the fewest pairs among the infinitesimally rigid sets (76); Eisenstein-33 has the most pairs (78) but the fewest triads (14); Peres-33 is intermediate (72~pairs, 16~triads) and is the only one that is not infinitesimally rigid.  This suggests that the triad/pair balance, rather than either quantity alone, determines infinitesimal rigidity.
\end{observation}

\begin{observation}[MUS landscape of the integer pool]\label{obs:mus-landscape}
Among 5{,}000 independent MUS (Minimal Unsatisfiable Subset) extractions from the 49-ray integer pool, exactly \emph{six} distinct minimal 31-sets are found.  These six sets collectively use 37 of the 49 rays, with a sharp norm-stratified structure: 13 rays with $\|v\|^2 \leq 3$ (the 3 axis directions, 6 face diagonals, and 4 body diagonals of the cube) appear in \emph{every} minimal set; 12 rays with $\|v\|^2 = 5$ appear in 5 of 6; 12 rays with $\|v\|^2 = 6$ appear in 4 of 6; and 12 rays with $\|v\|^2 = 9$ appear in \emph{none}.  The 13 core rays are precisely the $\{0, \pm 1\}$ alphabet---the minimal alphabet that is too small for KS-uncolorability on its own but forms the invariant skeleton of every minimal set.  All pairs of 31-sets share 25--26 rays (Jaccard similarity $\approx 0.71$), yet no two are connected by a single-ray swap (each pair differs by 5--6 rays).  CK-31 is one of the six.
\end{observation}

\begin{conjecture}\label{conj:optimal}
The Conway--Kochen set $\CK$ achieves the true minimum for KS sets in $\R^3$: no KS set with 30 or fewer rays exists.
\end{conjecture}

This conjecture is supported by: the rigidity of $\CK$ in~$\C^3$ (Observation~\ref{obs:rigidity}; confirmed independently by Trandafir and Cabello~\cite{TrandafirCabello2025}); the OCUS proof that no $\leq 30$-ray KS subset exists within the full 49-ray integer pool (Remark after Proposition~\ref{prop:31optimal}); the absence of any sub-31 KS set across all six algebraic islands (Observation~\ref{obs:islands}); universal merge saturation (Observation~\ref{obs:merge}); the MUS landscape analysis showing only six distinct minimal 31-sets, all sharing a 13-ray invariant core and exhibiting sharp norm stratification (Observation~\ref{obs:mus-landscape}); and the complete failure of 30{,}000+ random hypergraphs to achieve geometric realizability.  The graph universality result (Observation~\ref{obs:universality}) is a consistency check entailed by rigidity rather than independent evidence, but confirms that no alternative 31-vertex orthogonality graph emerges from any tested construction family.

\begin{observation}[Graph universality of $\CK$]\label{obs:universality}
We tested every algebraic construction in our survey that achieves the minimum of 31 rays---the integer alphabet $\{0,\pm 1,\pm 2\}$, the rational alphabet $\{0,\pm 1,\pm 2,\pm\tfrac{1}{2}\}$, the extended Peres alphabet $\{0,\pm 1,\pm\sqrt{2},\pm(1+\sqrt{2}),\pm 2\}$, the mixed modulus-2 alphabet $\{0,\pm 1,\pm 2,\pm\sqrt{2},\pm\omega,\pm\bar\omega\}$ combining three algebraic cancellation identities, and finite group orbits under the tetrahedral group $A_4$ and octahedral group $S_4$ acting on seed vectors in $\R^3$---and in every case the resulting minimal 31-ray KS set is graph-isomorphic (verified by VF2 with explicit vertex bijections) and hypergraph-isomorphic (bijections preserve triads): 71 orthogonality edges, 17 bases, degree sequence $(3^4, 4^{14}, 5^8, 6^3, 8^2)$. Despite arising from algebraically independent constructions, all roads to 31 lead to the same orthogonality graph.
\end{observation}

Since Trandafir and Cabello~\cite{TrandafirCabello2025} proved that $\CK$ is rigid in~$\C^3$ (unique up to unitary equivalence), the universality across algebraic constructions is \emph{expected}: any realization of a 31-vertex KS set must be unitarily equivalent to $\CK$, so all constructions that achieve 31 rays necessarily produce the same orthogonality graph.  The value of Observation~\ref{obs:universality} is therefore not as independent evidence for Conjecture~\ref{conj:optimal}, but as a \emph{consistency check}: it confirms that our diverse algebraic constructions all land on the unique rigid realization, and that no alternative 31-vertex orthogonality graph emerges from any tested family.  The 31-vertex graph appears to be the unique minimal realizable KS hypergraph in dimension~3.

\section{Limitations and Open Questions}\label{sec:limitations}

Our search is computational, not exhaustive over all possible number fields.  It is important to understand both its scope and its principled structure.

The relationship between alphabets, cancellation identities, and KS sets is tight: an alphabet is a candidate for producing a KS set \emph{only if} it carries a cancellation identity---an algebraic relation among its elements that makes dot products vanish exactly.  Without such an identity, the alphabet cannot produce exact orthogonalities among triples of rays, and without orthogonal triples there is no KS set.  The search for new KS-supporting alphabets is therefore equivalent to the search for new cancellation identities of sufficiently low norm.  Our survey is systematic within this framework: for each class of number fields, we identify the candidate cancellation identities from the ring structure and test whether they generate KS-uncolorable ray sets.

The coverage is thorough for low-degree extensions but becomes increasingly sparse at higher degree:

\begin{enumerate}
    \item \textbf{Quadratic fields (near-exhaustive)}: we tested real $\Q(\sqrt{d})$ for all non-square $d = 2, \ldots, 30$ with the basic alphabet $\{0, \pm 1, \pm\sqrt{d}\}$, and also the ring-of-integers generator $(1+\sqrt{d})/2$ for $d \equiv 1\pmod{4}$ ($d \in \{5, 13, 17, 21, 29\}$).  The golden ratio case ($d = 5$) demonstrates that the choice of generator within a field matters: $\varphi = (1+\sqrt{5})/2$ produces a KS set while $\sqrt{5}$ does not, because $\varphi$ participates in a cancellation identity ($\varphi^2 = \varphi + 1$) that $\sqrt{5}$ does not.
    \item \textbf{Complex quadratic fields (exhaustive for class number~1)}: we tested all nine imaginary quadratic fields with class number 1 (the Heegner numbers: $d = 1, 2, 3, 7, 11, 19, 43, 67, 163$), as well as $\Z[\sqrt{-d}]$ for $d = 2, 5, 6, 7, 10, 11, 13, 14, 15$ with enriched alphabets and Hermitian completion.  Higher class numbers (where the ring of integers is no longer a UFD---a \emph{unique factorization domain}, meaning a ring in which every nonzero non-unit element factors uniquely into irreducibles, up to order and units) remain largely untested; different factorization structures could in principle support cancellation identities not available in class-1 rings. A systematic survey of class-number $> 1$ fields is reported below (item~\ref{item:class-number-survey}).
    \item \textbf{Cyclotomic fields (good coverage)}: all roots of unity of order $n \leq 60$ were tested.  The $6 \mid n$ characterization (Theorem~\ref{thm:6n}) is proved algebraically for all four cases (sufficiency plus three necessity cases).
    \item \textbf{Cubic and higher extensions (sparse)}: we tested cubic extensions $\Q(\sqrt[3]{d})$ for $d = 2, 3, 4, 5, 6, 7$. The basic alphabets $\{0,\pm 1,\pm\sqrt[3]{d}\}$ are all colorable, but the extended alphabet $\{0,\pm 1,\pm\sqrt[3]{2},\pm\sqrt[3]{4}\}$ of $\Q(\sqrt[3]{2})$ produces an uncolorable set after cross-product completion (361 rays, 280 triads; raw pool is colorable).  Full characterization yields: minimum 58~vectors in 26~bases, 140~orthogonal pairs, 7~distinct degree types (range 3--12), and infinitesimal rigidity in~$\R^3$.  The cancellation mechanism is indirect modulus-2: $\sqrt[3]{2} \cdot \sqrt[3]{4} = 2$, producing the same $1 + 1 - 2 = 0$ identity through the \emph{product} of two generators rather than the squared modulus of one.  The BPQS cost ($13 \times 13 = 169$) is the highest of any characterized island, and the orthogonality graph is not isomorphic to any of the six known islands (different vertex count).  This confirms that cubic fields can support KS sets through indirect modulus-2 cancellation, but at sharply increased cost (58~vectors vs.\ 31--52 for quadratic islands).  Quartic extensions were initially untested; three quartic cases are now covered (item~\ref{item:class-number-survey}).  Quintic and higher-degree extensions remain unexplored; these could harbor cancellation identities qualitatively different from the quadratic modulus-2 pattern.  This is the principal gap in our coverage.
    \item \textbf{Exotic cancellation candidates}: cubic Pisot-number alphabets (e.g., the plastic ratio $\rho$, root of $x^3 = x+1$, with irreducibly cubic cancellation $\rho \cdot \rho^2 - \rho - 1 = 0$) and multi-generator alphabets with simultaneous conjugate cancellations remain largely untested.  Preliminary tests of Pisot-number families found no KS sets; scripts are included in the repository.
    \item \label{item:class-number-survey}\textbf{Class-number $> 1$ survey (systematic)}: we tested 26 imaginary quadratic fields $\Q(\sqrt{-d})$ with class number $h \in \{2, 3, 4\}$, covering all such fields with $d \leq 91$ (for $h = 2$), $d \leq 83$ (for $h = 3$), and $d \leq 57$ (for $h = 4$). For each field we used the correct ring-of-integers generator: $\sqrt{-d}$ when $d \not\equiv 3 \pmod{4}$, and $(1 + \sqrt{-d})/2$ when $d \equiv 3 \pmod{4}$. Both a basic alphabet $\{0, \pm 1, \pm\alpha, \pm\bar\alpha\}$ and an enriched alphabet (adding products and sums of generators, up to 15 elements) were tested. Of the 20 genuinely new fields, \emph{all are colorable at both alphabet levels} except $d = 15$ ($h = 2$, $|\alpha|^2 = 4$, enriched alphabet uncolorable).  A deep dive into $d = 15$ reveals that the enriched 15-element alphabet generates 1{,}537 rays with 254 triads; greedy minimization over 50 trials yields a minimum of~85 vectors---nearly twice the Golden island minimum---confirming that enriched alphabets can bootstrap KS-uncolorability from generators above the norm-2 threshold, but at sharply increased cost.  No norm-2 products exist in the $d = 15$ product set; the cancellations involve terms of squared modulus 4, 6, and~16. The previously tested fields $d \in \{5, 6\}$ also remain uncolorable with enriched alphabets (as expected from their low generator norms $|\alpha|^2 = 5$ and~6 respectively), while $d \in \{10, 13, 14\}$ remain colorable. All fields with $h = 3$ and $h = 4$ are colorable.

    The pattern is governed by generator norm, not class number: fields with $|\alpha|^2 \leq 6$ can support KS sets (though not all do), while fields with $|\alpha|^2 \geq 8$ are uniformly colorable across all tested alphabets. This is consistent with the modulus-2 cancellation thesis---the enriched alphabets of low-norm fields contain elements whose products reach the critical $|\cdot|^2 = 2$ threshold, while high-norm generators produce squared moduli too large for the necessary dot-product cancellations. The class-number distinction (UFD vs.\ non-UFD) does not independently predict KS-uncolorability.

    Additionally, we tested three quartic fields $\Q(d^{1/4})$ for $d = 2, 3, 5$ with extended alphabets including all powers $\{d^{1/4}, d^{1/2}, d^{3/4}\}$: only $d = 2$ is uncolorable (minimum 33 vectors), consistent with the $\sqrt{2}$ modulus-2 mechanism. We also surveyed 28 two-generator alphabets formed from pairs of $\{\sqrt{2}, \sqrt{3}, \sqrt{5}, \varphi, \omega, \sqrt{-2}, i, (1{+}\sqrt{-7})/2\}$: 22 of 28 are uncolorable.  Among the 17 pairs achieving minimum~33, VF2 isomorphism testing reveals exactly two graph classes: 13~pairs containing a modulus-2 generator ($\sqrt{2}$ or~$\sqrt{-2}$) produce the Peres graph (degree distribution $\{4{:}30, 8{:}3\}$, 72~pairs, 16~triads), while 4~pairs containing $\omega$ produce the Eisenstein graph ($\{4{:}18, 5{:}12, 8{:}3\}$, 78~pairs, 14~triads).  The 3~pairs achieving minimum~43 (all containing the Heegner-7 generator) are mutually graph-isomorphic.  This confirms that the cancellation identity, not the ambient field, determines the orthogonality graph.  Six pairs are colorable: those pairing $i$, $\varphi$, $\sqrt{3}$, or $\sqrt{5}$ without any modulus-2 or phase-cancellation generator.
    \item \textbf{Non-commutative division algebras}: quaternionic and octonionic alphabets in $\mathbb{H}^3$ and $\mathbb{O}^3$ produce KS-uncolorable pools when they contain the modulus-2 identity, but greedy minimization always converges to 31~vectors with the CK-31 graph.  Non-commutativity and non-associativity destroy orthogonalities that would exist over~$\C$, and this loss outweighs the additional cancellation identities.  The modulus-2 boundary appears universal across all four normed division algebras.\label{rem:division-algebras}
    \item \textbf{Completion termination}: cross-product completion terminates empirically for all tested alphabets, but we do not prove termination in general.
    \item \textbf{Classification scope}: the ``island'' classification is a classification of \emph{coordinate realizations}, not of abstract KS graphs. The intrinsic object is the orthogonality hypergraph; a given hypergraph may admit realizations in multiple number fields (as demonstrated by the Peres/$\Z[\sqrt{-2}]$ isomorphism). Conversely, different coordinatizations of the same abstract graph may have different operational properties (BPQS efficiency, CSW advantage). Our contribution is that the algebraic substrate---the cancellation identity available in the coordinate ring---determines which graphs can be realized, and hence which operational features are accessible.
    \item \textbf{Transcendental coordinates}: our methods require algebraic coordinates (for exact orthogonality testing); KS sets with transcendental coordinates are not excluded.
    \item \textbf{The $6 \mid n$ theorem}: Theorem~\ref{thm:6n} is proved algebraically for all four cases.  It is the only fully proved result in this paper; all other claims about minimality and optimality are computational.  The proof was constructed and checked using LLM assistance; independent human verification is invited, particularly for the necessity cases (Cases~2 and~3), which required multiple correction rounds during development.
    \item \textbf{Minimization is now certified}: the OCUS (Optimal Constrained Unsatisfiable Subset) procedure exhaustively certifies that each pool minimum in Observation~\ref{obs:islands} is the exact minimum within its pool. For the first five pools (Integer, Peres, Eisenstein, $\Z[\sqrt{-2}]$, Heegner-7), certification completes in under 4\,s each; the Golden pool required 18{,}783 CEGAR iterations (303\,s) at $n = 51$. These are \emph{proofs within each finite pool}, not heuristic upper bounds. However, they do not exclude the possibility that a smaller KS set could exist using coordinates outside these pools---e.g., from an untested number field.
    \item \textbf{Lower bound improvement}: our work does not improve the lower bound of 24; it constrains where a sub-31 construction could come from.  These constraints---the cancellation classification (Observation~\ref{obs:two-element}), density monotonicity, and the discrete 13$\to$49 ray jump---are directly applicable as filters for the realizability-based programs of Li, Bright, and Ganesh~\cite{LiBrightGanesh2024} and Kirchweger et al.~\cite{KirchwegarPeitlSzeider2023}.
    \item \textbf{Operational significance}: our connection to Cabello's work (Section~\ref{sec:cabello}) is primarily observational---we identify the Eisenstein island with his simplest KS set, but do not prove that other islands cannot arise as engines of perfect quantum strategies. A systematic study of which Bell scenarios' perfect strategies correspond to which algebraic islands would deepen this connection.
    \item \textbf{Pareto structure of the trade-off}: A \emph{Pareto frontier} is the set of options for which no other option is simultaneously better on every criterion; any improvement on one axis requires a sacrifice on another.  The three-way trade-off among the islands (ray economy, BPQS cost, and CSW pool advantage) exhibits a clean Pareto structure.  Three islands lie on the frontier: $\CK$ (fewest rays, 31), Eisenstein (lowest BPQS product, $5 \times 9 = 45$), and Heegner-7 (strongest CSW advantage, $\vartheta/\alpha = 1.118$).  The Peres and $\Z[\sqrt{-2}]$ islands are Pareto-dominated: they match Eisenstein's ray count (33) but have higher BPQS cost ($7 \times 9 = 63$) and no CSW advantage ($\vartheta = \alpha$).  The Golden island is dominated on all three axes.  This confirms that the choice of algebraic substrate has genuine operational consequences and that no single island is ``best'' in all senses.
    \item \textbf{CK-31 and perfect quantum strategies}: A systematic B-KS search over 200 greedy trials shows that every Bell scenario achievable by $\CK$ (minimum $8 \times 9 = 72$) is also achievable at equal or lower cost by the Eisenstein island (minimum $5 \times 9 = 45$) and the Peres island (minimum $7 \times 9 = 63$).  $\CK$ requires at least 8~Alice inputs, versus 5 for Eisenstein and 7 for Peres.  Thus $\CK$, despite having the fewest rays, is strictly dominated as a BPQS engine: the 33-vector islands generate simpler Bell scenarios. The operational advantage of $\CK$ lies in its unique contextual tightness ($\eta = 1.000$) and CSW sandwich ($\alpha = 11 < \vartheta = 11.71 < \alpha^* = 12$), not in BPQS economy.
\end{enumerate}

\section{Conclusion}\label{sec:conclusion}

The central empirical finding of this paper is that, among all tested coordinate alphabets, KS-uncolorability in dimension~3 requires one of exactly two cancellation mechanisms: \emph{modulus-2 cancellation} (an algebraic identity producing the value~2 from products of ring elements, as in $1+1=2$, $(\sqrt{2})^2=2$, or $\alpha\bar\alpha=2$) or \emph{phase cancellation} (a root-of-unity sum vanishing exactly, as in $1+\omega+\omega^2=0$).  Among all quadratic fields, cyclotomic fields, and selected extensions tested, KS sets exist only when the coordinate ring supports one of these two mechanisms; alphabets where the smallest available cancellation involves norm~$\geq 3$ produce orthogonal triples but not KS-uncolorability.  This two-mechanism constraint explains why constructions cluster into at least six discrete algebraic islands among the fields tested (a seventh, cubic island at 58~vectors is characterized in Section~\ref{sec:limitations}).  The golden ratio island ($|\varphi|^2 \approx 2.618$) does not violate this boundary: cross-product completion introduces modulus-2-type cancellations into the effective coordinate set, confirming that the relevant criterion is the cancellation structure of the completed alphabet, not the generator's own norm.

A systematic survey of 26 imaginary quadratic fields with class number 2, 3, and 4 found no new islands: all fields with generator norm $|\alpha|^2 \geq 8$ are colorable, confirming that the operative constraint is generator norm rather than factorization structure (Section~\ref{sec:limitations}, item~\ref{item:class-number-survey}).

This structural perspective yields both new constructions and new connections:

\begin{enumerate}
    \item \textbf{New KS configurations.} The Heegner-7 ring $\Z[(1+\sqrt{-7})/2]$ produces a genuinely new orthogonality graph (43 vectors, 23 bases), and the golden ratio field $\Q(\varphi)$ yields a KS set (52 vectors) that is invisible to standard alphabet searches and revealed only by cross-product completion. Neither has been previously reported. The complex quadratic $\Z[\sqrt{-2}]$ provides a new algebraic realization of the Peres-type modulus-2 structure, graph-isomorphic to the original (Proposition~\ref{prop:isomorphism}).

    \item \textbf{Operational consequences.} The algebraic substrate determines the complexity of the resulting quantum information tasks. Using SAT-based bipartite KS-uncolorability, we verify and extend the BPQS input counts of Trandafir and Cabello (Table~\ref{tab:bpqs}): the Eisenstein island yields the simplest Bell scenario ($5 \times 9$), while CK-31 is costlier ($\leq 8 \times 9$) despite having fewer rays. CSW analysis shows that auxiliary rays amplify contextual advantage, and that no single island dominates across all operational measures---ray economy, context economy, and CSW strength are in tension.

    \item \textbf{Supporting evidence.} The $6 \mid n$ characterization for cyclotomic KS-uncolorability (Theorem~\ref{thm:6n}) connects to the $\Z[1/6]$ minimality result of Cortez, Morales, and Reyes, with the Eisenstein island achieving exactly $N(S) = 6$.  An OCUS procedure exhaustively proves that no $\leq 30$-ray KS subset exists within the 49-ray integer pool.  All six minimal KS sets are merge-saturated (Observation~\ref{obs:merge}), four of six are rigid in~$\C^3$ (Observation~\ref{obs:rigidity}), the integer pool contains exactly six distinct minimal 31-sets with a norm-stratified 13-ray invariant core (Observation~\ref{obs:mus-landscape}), and three structurally distinct 33-vector KS sets exist with pairwise non-isomorphic graphs (Observation~\ref{obs:three-33}).  Graph universality (Observation~\ref{obs:universality}) confirms that every tested construction achieving 31~rays produces the same orthogonality graph.
\end{enumerate}

Whether the two-mechanism constraint extends to all number fields, or whether higher-degree extensions harbor additional cancellation mechanisms and hence additional islands, remains the principal open question. The identification of Cabello's ``simplest KS set'' with the Eisenstein island suggests that the algebraic perspective may complement the ongoing search for optimal constructions. More concretely, the constraints established here---the two-element cancellation classification, density monotonicity, and the discrete 13$\to$49 ray jump---are directly applicable to the realizability-based lower-bound program of Li, Bright, and Ganesh~\cite{LiBrightGanesh2024}: algebraic pruning of candidate hypergraphs could substantially narrow the search space that their SMT solver must explore, bringing the constructive upper bound (31) and the realizability lower bound (24) closer together.

The two-mechanism thesis is falsifiable: a KS-uncolorable set from an alphabet whose product set admits no modulus-2 or phase cancellation identity would refute it. The most promising candidates were cubic Pisot-number alphabets (where the cancellation identity is the minimal polynomial itself), class-number $>1$ imaginary quadratic fields (where non-unique factorization provides alternative algebraic relations), and multi-generator alphabets with cross-cancellation mechanisms. A systematic survey of 26 class-number $> 1$ fields (Section~\ref{sec:limitations}, item~\ref{item:class-number-survey}) found no counterexamples: all tested non-UFD rings are colorable when $|\alpha|^2 \geq 8$, confirming that the operative constraint is generator norm rather than factorization structure. The space of higher-degree extensions remains the principal open frontier.

\subsection*{Note added in revision}

\begin{sloppypar}
After this paper first appeared, Gunji et al.~\cite{GunjiOhzawa2026} showed that orthomodular (non-distributive) lattice structure arises canonically as a left adjoint from classical Boolean contexts via categorical pushouts, and that the failure of such a pushout to remain Boolean is equivalent to the absence of global sections in the Abramsky--Brandenburger sheaf-theoretic framework.  Their result provides a category-theoretic explanation for \emph{why} contextuality forces non-Boolean structure, complementing the present paper's computational classification of \emph{which} coordinate algebras support KS-uncolorability.  It is natural to ask whether the algebraic islands identified here correspond to settings whose orthogonality hypergraphs, viewed categorically as contexts, yield non-trivial (i.e., non-Boolean) pushouts in the sense of~\cite{GunjiOhzawa2026}; we leave a precise formulation and proof of this correspondence to future work.
\end{sloppypar}

\subsection*{Note added in second revision}

The vanishing-sum enumeration in Table~\ref{tab:two-element} has been corrected to include the previously missing zero-sum $1 - x - \bar{x} = 0$ (constraint $\mathrm{Tr}(x) = 1$), which arises from the sign pattern $(+1, -1, -1)$ on the triple $\{1, x, \bar{x}\}$.  This pattern produces all-nonzero orthogonal triads for imaginary quadratic generators with $d \equiv 1 \pmod{4}$, but does not lead to KS-uncolorability: the cross-product obstruction (no type-B triad can form) is now proved for \emph{all} integer generator norms $N \geq 3$ via SMT verification~\cite{deMouraBjorner2008}, upgrading the previous spot-checks to a universal certificate.  No computational results or conclusions are affected; the correction tightens the enumeration and strengthens the proof.

\subsection*{Reproducibility}

\begin{sloppypar}
All code and data are available at
\url{https://github.com/michaelkernaghan/contextuality}.
The repository includes explicit ray coordinates and triad/basis lists for all islands, including the 43-vector Heegner-7 and 52-vector golden-ratio sets, as well as verifier scripts that reconstruct each KS set from its alphabet and confirm UNSAT via SAT solving. The script \texttt{ks\_z3\_verify.py} provides SMT-certified proofs (Z3~\cite{deMouraBjorner2008}) that the vanishing-sum enumeration is complete and that the cross-product obstruction holds for all quadratic integer generators with $N \geq 3$. The implementation uses Python~3.11 with PySAT~1.8 (Glucose4), Z3~4.15, NumPy, SciPy (L-BFGS-B and HiGHS~LP), and CVXPY (SCS~SDP, tolerance~$10^{-9}$).
\end{sloppypar}

\begin{sloppypar}
\noindent The script-to-table correspondence is:
\begin{itemize}\setlength{\itemsep}{0pt}\setlength{\parskip}{0pt}
\item \texttt{ks\_islands.py} $\to$ Tables~\ref{tab:alphabets}--\ref{tab:closure}
\item \texttt{ks\_complex.py} $\to$ Table~\ref{tab:roots}
\item \texttt{ks\_new\_islands.py}, \texttt{ks\_sat.py} $\to$ Table~\ref{tab:complex-quadratic}, 31-optimality check
\item \texttt{ks\_new\_island\_analysis.py} $\to$ Table~\ref{tab:heegner}
\item \texttt{ks\_csw.py}, \texttt{ks\_csw\_extended.py}, \texttt{ks\_graph\_analysis.py} $\to$ Tables~\ref{tab:csw-min}--\ref{tab:csw-pool},~\ref{tab:spectral}
\item \texttt{ks\_critical\_bases.py} $\to$ Table~\ref{tab:critical-bases}
\item \texttt{ks\_bpqs\_sat.py} $\to$ Table~\ref{tab:bpqs}
\item \texttt{ks\_trig.py} $\to$ Remark~\ref{rem:trig}
\item \texttt{ks\_graph\_analysis.py} $\to$ Prop.~\ref{prop:isomorphism}
\item \texttt{ks\_group\_isomorphism.py} $\to$ Obs.~\ref{obs:universality}
\item \texttt{ks\_explore\_new*.py} $\to$ cubic/mixed/orbit explorations
\item \texttt{ks\_ocus\_all\_pools.py} $\to$ OCUS certification of pool minima (Obs.~\ref{obs:islands})
\item \texttt{ks\_maxsat\_optimal.py} $\to$ Integer-pool OCUS (Remark after Prop.~\ref{prop:31optimal})
\item \texttt{ks\_peres\_merge.py} $\to$ Obs.~\ref{obs:merge}
\item \texttt{ks\_literature\_connections.py} $\to$ SI-C closure, $N(S)$ invariant
\item \texttt{ks\_mus\_landscape.py} $\to$ Obs.~\ref{obs:mus-landscape}
\item \texttt{ks\_rigidity.py}, \texttt{ks\_rigidity\_finite.py}, \texttt{ks\_rigidity\_global\_search.py} $\to$ Obs.~\ref{obs:rigidity}
\item \texttt{ks\_d15\_deep\_dive.py} $\to$ Limitations, item~\ref{item:class-number-survey}
\item \texttt{ks\_two\_gen\_isomorphism.py} $\to$ Limitations, item~\ref{item:class-number-survey} (VF2 check)
\item \texttt{ks\_bpqs\_unique.py} $\to$ Limitations, CK-31 BPQS uniqueness analysis
\item \texttt{ks\_cubic\_characterize.py} $\to$ Limitations, item~4 (cubic island characterization)
\item \texttt{ks\_ck33\_identify.py} $\to$ Obs.~\ref{obs:three-33}
\item \texttt{ks\_quaternion.py}, \texttt{ks\_octonion.py} $\to$ Limitations, item~\ref{rem:division-algebras}
\item \texttt{ks\_6n\_proof.py}, \texttt{ks\_6n\_proof\_theorem.py} $\to$ Thm.~\ref{thm:6n}
\item \texttt{ks\_z3\_verify.py} $\to$ SMT certification of Table~\ref{tab:two-element} completeness and cross-product obstruction (Obs.~\ref{obs:two-element})
\item \texttt{ks\_verify\_galois.py} $\to$ SymPy/SAT verification of Galois letter proofs (companion paper)
\end{itemize}
\end{sloppypar}
All randomized experiments use \texttt{random.seed(42)} for~reproducibility.

\subsection*{Verification methodology}

This paper was developed using an LLM-assisted workflow in which all computations, proofs, and manuscript text were generated by a large language model (Claude, Anthropic) under the author's direction.  To mitigate known failure modes of LLM-generated mathematics---particularly silent fabrication and pressure-induced confabulation~\cite{Schwartz2026}---three verification layers were applied: (1)~all computational claims are independently reproducible via the scripts listed above with deterministic seeding; (2)~the algebraic proof of Theorem~\ref{thm:6n} was subjected to six rounds of adversarial peer review using a competing LLM (GPT, OpenAI), with substantive errors found and corrected in rounds~2 and~5; (3)~all claims about published results were verified by direct reading of the cited papers.  The author invites independent verification of Theorem~\ref{thm:6n}, which is the only fully proved result; all other claims are computational and reproducible.

\subsection*{Acknowledgments}

The author thanks Ad\'an Cabello for pointing out the connection between the Eisenstein island and his ``simplest KS set'' construction.  Claude (Anthropic) performed all computational work including Python script development, SAT-based verification, rigidity analysis, algebraic proof construction for Theorem~\ref{thm:6n}, and manuscript preparation; cross-verification was performed using GPT (OpenAI) through structured adversarial peer review.  The author directed the research program, selected problems, validated results against the published literature, and accepts full responsibility for the scientific content and integrity of this paper.

\begin{thebibliography}{99}

\bibitem{KochenSpecker1967}
S.~Kochen and E.~P.~Specker,
``The problem of hidden variables in quantum mechanics,''
\textit{J. Math. Mech.} \textbf{17}, 59--87 (1967).

\bibitem{Peres1991}
A.~Peres,
``Two simple proofs of the Kochen--Specker theorem,''
\textit{J. Phys. A: Math. Gen.} \textbf{24}, L175--L178 (1991).

\bibitem{Peres1993}
A.~Peres,
\textit{Quantum Theory: Concepts and Methods}
(Kluwer Academic Publishers, 1993).

\bibitem{Kernaghan1994}
M.~Kernaghan,
``Bell--Kochen--Specker theorem for 20 vectors,''
\textit{J. Phys. A: Math. Gen.} \textbf{27}, L829--L830 (1994).

\bibitem{KernaghanPeres1995}
M.~Kernaghan and A.~Peres,
``Kochen--Specker theorem for eight-dimensional space,''
\textit{Phys. Lett. A} \textbf{198}, 1--5 (1995).

\bibitem{CabelloEstebaranzGarcia1996}
A.~Cabello, J.~M.~Estebaranz, and G.~Garc\'{\i}a-Alcaine,
``Bell--Kochen--Specker theorem: A proof with 18 vectors,''
\textit{Phys. Lett. A} \textbf{212}, 183--187 (1996).

\bibitem{Pavicic2004}
M.~Pavicic, J.-P.~Merlet, B.~McKay, and N.~D.~Megill,
``Kochen--Specker vectors,''
\textit{J. Phys. A: Math. Gen.} \textbf{38}, 1577--1592 (2005).
arXiv:quant-ph/0409014.

\bibitem{Pavicic2019}
M.~Pavicic, N.~D.~Megill, B.~C.~Waegell, and P.~K.~Aravind,
``Automated generation of Kochen--Specker sets,''
\textit{Sci. Rep.} \textbf{9}, 6765 (2019).
arXiv:1905.01009.

\bibitem{CabelloSeveriniWinter2014}
A.~Cabello, S.~Severini, and A.~Winter,
``Graph-theoretic approach to quantum correlations,''
\textit{Phys. Rev. Lett.} \textbf{112}, 040401 (2014).

\bibitem{UijlenWesterbaan2016}
S.~Uijlen and B.~Westerbaan,
``A Kochen--Specker system has at least 22 vectors,''
\textit{New Generation Computing} \textbf{34}, 3--23 (2016).

\bibitem{AudemardSimon2018}
G.~Audemard and L.~Simon,
``On the Glucose SAT solver,''
\textit{Int. J. Artif. Intell. Tools} \textbf{27}, 1840001 (2018).

\bibitem{KunjwalSpekkens2018}
R.~Kunjwal and R.~W.~Spekkens,
``Beyond the Cabello--Severini--Winter framework: Making sense of contextuality without sharpness of measurements,''
\textit{Phys. Rev. A} \textbf{97}, 052110 (2018).

\bibitem{CortezMoralesReyes2022}
I.~Cortez, C.~Morales, and M.~Reyes,
``Minimal ring extensions of the integers exhibiting Kochen--Specker contextuality,''
arXiv:2211.13216v4 [math.NT] (2025, preprint).

\bibitem{LiBrightGanesh2024}
Z.~Li, C.~Bright, and V.~Ganesh,
``A {SAT} solver + computer algebra attack on the minimum {Kochen}--{Specker} problem,''
in \textit{Proceedings of the Thirty-Third International Joint Conference on Artificial Intelligence, {IJCAI-24}},
ed.\ K.~Larson (International Joint Conferences on Artificial Intelligence Organization, 2024),
pp.~1898--1906.
\textsc{doi}: 10.24963/ijcai.2024/210.

\bibitem{Cabello2024bipartite}
A.~Cabello,
``Simplest bipartite perfect quantum strategies,''
\textit{Phys.\ Rev.\ Lett.} \textbf{134}, 010201 (2025).
arXiv:2311.17735.

\bibitem{Cabello2024equivalences}
Y.~Liu, H.~Y.~Chung, E.~Z.~Cruzeiro, J.~R.~Gonzales-Ureta, R.~Ramanathan, and A.~Cabello,
``Equivalence between face nonsignaling correlations, full nonlocality, all-versus-nothing proofs, and pseudotelepathy,''
\textit{Phys.\ Rev.\ Research} \textbf{6}, L042035 (2024).
arXiv:2310.10600.

\bibitem{Cabello2025simplest}
A.~Cabello,
``Simplest Kochen--Specker set,''
\textit{Phys.\ Rev.\ Lett.} (in press, 2025).
arXiv:2508.07335.

\bibitem{TrandafirCabello2025}
S.~Trandafir and A.~Cabello,
``Two fundamental solutions to the rigid Kochen--Specker set problem and the solution to the minimal Kochen--Specker set problem under one assumption,''
\textit{Phys.\ Rev.\ A} \textbf{111}, 052204 (2025).
arXiv:2501.11640.

\bibitem{TrandafirCabello2025optimal}
S.~Trandafir and A.~Cabello,
``Optimal conversion of Kochen--Specker sets into bipartite perfect quantum strategies,''
\textit{Phys.\ Rev.\ A} \textbf{111}, 022408 (2025).
arXiv:2410.17470.

\bibitem{AbramskyBarbosaMansfield2017}
S.~Abramsky, R.~S.~Barbosa, and S.~Mansfield,
``The contextual fraction as a measure of contextuality,''
\textit{Phys.\ Rev.\ Lett.} \textbf{119}, 050504 (2017).
arXiv:1705.07918.

\bibitem{Gleason1957}
A.~M.~Gleason,
``Measures on the closed subspaces of a Hilbert space,''
\textit{J. Math. Mech.} \textbf{6}, 885--893 (1957).

\bibitem{RajanVisser2019}
D.~Rajan and M.~Visser,
``Kochen--Specker theorem revisited,''
arXiv:1708.01380 [quant-ph] (2019).

\bibitem{Haemers1995}
W.~H.~Haemers,
``Interlacing eigenvalues and graphs,''
\textit{Linear Algebra Appl.} \textbf{226--228}, 593--616 (1995).

\bibitem{Knuth1994}
D.~E.~Knuth,
``The sandwich theorem,''
\textit{Electron.\ J.\ Combin.} \textbf{1}, Article A1 (1994).

\bibitem{GouldAravind2010}
P.~K.~Gould and P.~K.~Aravind,
``Isomorphism between the Peres and Penrose proofs of the BKS theorem in three dimensions,''
\textit{Found.\ Phys.} \textbf{40}, 1096--1101 (2010).

\bibitem{KirchwegarPeitlSzeider2023}
M.~Kirchweger, A.~Peitl, and S.~Szeider,
``Co-certificate learning with SAT modulo symmetries,''
in \textit{Proceedings of AAAI-23}, 4119--4126 (2023).

\bibitem{Budroni2022}
C.~Budroni, A.~Cabello, O.~G\"uhne, M.~Kleinmann, and J.-\AA.~Larsson,
``Kochen--Specker contextuality,''
\textit{Rev.\ Mod.\ Phys.} \textbf{94}, 045007 (2022).

\bibitem{Cordella2004VF2}
L.~P.~Cordella, P.~Foggia, C.~Sansone, and M.~Vento,
``A (sub)graph isomorphism algorithm for matching large graphs,''
\textit{IEEE Trans.\ Pattern Anal.\ Mach.\ Intell.} \textbf{26}, 1367 (2004).

\bibitem{LamLeung2000}
T.~Y.~Lam and K.~H.~Leung,
``On vanishing sums of roots of unity,''
\textit{J.\ Algebra} \textbf{224}, 91 (2000).

\bibitem{BravyiGossetKoenig2018}
S.~Bravyi, D.~Gosset, and R.~K\"onig,
``Quantum advantage with shallow circuits,''
\textit{Science} \textbf{362}, 308 (2018).

\bibitem{Schwartz2026}
M.~D.~Schwartz,
``Vibe physics: the AI grad student,''
Anthropic Research (2026).
Available at \url{https://www.anthropic.com/research/vibe-physics}.
See also arXiv:2601.02484 for the resulting physics paper.

\bibitem{GunjiOhzawa2026}
Y.-P.~Gunji, Y.~Ohzawa, Y.~Tokuyama, Y.~Huang, and K.~Nakamura,
``Contextuality as a left adjoint: a categorical generation of orthomodular structure,''
arXiv:2603.22353 [quant-ph] (2026, preprint).

\bibitem{deMouraBjorner2008}
L.~de~Moura and N.~Bj{\o}rner,
``Z3: An efficient SMT solver,''
in \textit{TACAS 2008}, LNCS \textbf{4963}, pp.~337--340 (Springer, 2008).

\end{thebibliography}

\appendix

\section{Representative minimal KS sets for new islands}\label{app:rays}

For reproducibility, we describe the minimal KS sets found for the two genuinely new orthogonality graph types.

\textbf{Heegner-7 island} (43 vectors, 23 bases). Let $\alpha = (1+\sqrt{-7})/2$. All ray coordinates are drawn from $\{0, \pm 1, \pm\alpha, \pm\bar\alpha\}$, with cross-product completion adding elements of $\Z[\alpha]$ (integer linear combinations of $1$ and $\alpha$). The 145-ray pool is generated by enumerating all projectively distinct nonzero vectors $(a, b, c) \in \{0, \pm 1, \pm\alpha, \pm\bar\alpha\}^3$ under Hermitian orthogonality, then closing under cross products. The 43-ray minimal set is obtained by SAT-based greedy minimization (1000 trials); all single-ray removals from this set restore colorability. The 23 bases (orthogonal triples) and the degree distribution $\{4{:}16,\; 5{:}20,\; 6{:}4,\; 8{:}2,\; 10{:}1\}$ distinguish this hypergraph from all other known 3D KS configurations.

\textbf{Golden island} (52 vectors, 25 bases). Let $\varphi = (1+\sqrt{5})/2$. The raw alphabet $\{0, \pm 1, \pm\varphi\}$ generates 49 rays that are colorable; cross-product completion adds rays with coordinates in $\Z[\varphi] = \{a + b\varphi : a, b \in \Z\}$, expanding the pool to 205 rays. The completed pool is KS-uncolorable, with the smallest subset found containing 52 vectors in 25 bases. Completion-generated coordinates include $\varphi^2 = \varphi + 1$ and $2\varphi - 1 = \sqrt{5}$, both elements of $\Z[\varphi]$.

Complete ray coordinates (in exact algebraic form), basis hyperedge lists, and SAT verification scripts for all six islands are provided in the accompanying code repository.

\end{document}
